\documentclass[12pt,a4paper]{article}
\usepackage[top=10mm,left=15mm,right=15mm,bottom=15mm]{geometry}
\usepackage{amsmath,amsthm,amssymb,amsfonts}
\usepackage{booktabs,array}
\usepackage{enumitem}
\usepackage{xcolor}
\usepackage{tocloft}
% Compact TOC
\setlength{\cftbeforesecskip}{2pt}
\setlength{\cftbeforesubsecskip}{0pt}
\renewcommand{\cftsecleader}{\cftdotfill{\cftdotsep}}
\usepackage{hyperref}
\hypersetup{colorlinks=true,linkcolor=blue!70!black,urlcolor=blue!80!black,citecolor=blue!70!black}

% Theorem environments
\theoremstyle{definition}
\newtheorem{definition}{Definition}[section]
\newtheorem{example}{Example}[section]
\newtheorem{exercise}{Exercise}[section]
\newtheorem{theorem}{Theorem}[section]
\newtheorem{lemma}{Lemma}[section]
\newtheorem{corollary}{Corollary}[section]
\newtheorem{proposition}{Proposition}[section]
\newtheorem{property}{Property}[section]
\newtheorem*{remark}{Remark}

% Non-italic proof label
\makeatletter
\renewenvironment{proof}[1][\proofname]{\par
  \pushQED{\qed}%
  \normalfont \topsep6\p@\@plus6\p@\relax
  \trivlist
  \item[\hskip\labelsep\bfseries #1\@addpunct{.}]\ignorespaces
}{%
  \popQED\endtrivlist\@endpefalse
}
\makeatother

% Custom commands (matching clm.tex)
\newcommand\expc{\mathsf{E}}
\newcommand\prb{\mathsf{P}}
\DeclareMathOperator\var{var}
\DeclareMathOperator\cov{cov}
\DeclareMathOperator\corr{corr}
\DeclareMathOperator\plim{plim}
\DeclareMathOperator\tr{tr}
\DeclareMathOperator\rank{rank}
\DeclareMathOperator\diag{diag}
\newcommand{\ds}{\displaystyle}
\newcommand{\ie}{\;\Longrightarrow\;}
\newcommand{\ifff}{\;\Longleftrightarrow\;}
\newcommand{\SST}{\mathrm{SST}}
\newcommand{\SSR}{\mathrm{SSR}}
\newcommand{\SSE}{\mathrm{SSE}}
\newcommand{\MSR}{\mathrm{MSR}}
\newcommand{\MSE}{\mathrm{MSE}}
\newcommand{\SE}{\mathrm{SE}}
\newcommand{\bX}{\mathbf{X}}
\newcommand{\bY}{\mathbf{Y}}
\newcommand{\by}{\mathbf{y}}
\newcommand{\bx}{\mathbf{x}}
\def\bb{\boldsymbol{\beta}}
\newcommand{\be}{\boldsymbol{\varepsilon}}
\newcommand{\bI}{\mathbf{I}}
\newcommand{\bH}{\mathbf{H}}
\newcommand{\bM}{\mathbf{M}}
\newcommand{\bR}{\mathbf{R}}
\newcommand{\br}{\mathbf{r}}
\def\bc{\mathbf{c}}
\newcommand{\bone}{\mathbf{1}}
\newcommand{\bzero}{\mathbf{0}}
\newcommand{\bz}{\mathbf{z}}
\newcommand{\bw}{\mathbf{w}}
\newcommand{\bZ}{\mathbf{Z}}
\newcommand{\bA}{\mathbf{A}}
\newcommand{\bB}{\mathbf{B}}
\newcommand{\bC}{\mathbf{C}}
\newcommand{\bD}{\mathbf{D}}
\newcommand{\bP}{\mathbf{P}}
\newcommand{\bg}{\mathbf{g}}
\newcommand{\bh}{\mathbf{h}}
\newcommand{\bU}{\mathbf{U}}
\newcommand{\bLambda}{\boldsymbol{\Lambda}}
\newcommand{\ba}{\mathbf{a}}
\newcommand{\bv}{\mathbf{v}}
\newcommand{\bu}{\mathbf{u}}
\newcommand{\Real}{\mathbb{R}}
\newcommand{\calC}{\mathcal{C}}
\newcommand{\calN}{\mathcal{N}}
% Additional commands for modern theory
\newcommand{\bSigma}{\boldsymbol{\Sigma}}
\newcommand{\hbSigma}{\hat{\boldsymbol{\Sigma}}}
\newcommand{\hbb}{\hat{\boldsymbol{\beta}}}
\newcommand{\ip}[2]{\langle #1,\, #2 \rangle}
\newcommand{\norm}[1]{\|#1\|}
\newcommand{\opnorm}[1]{\|#1\|_{\mathrm{op}}}
\newcommand{\fnorm}[1]{\|#1\|_{\mathrm{F}}}
\DeclareMathOperator{\sgn}{sgn}
\DeclareMathOperator*{\esssup}{ess\,sup}
\newcommand{\dd}{\mathrm{d}}

\usepackage{parskip}
\usepackage{etoolbox}
\AtBeginEnvironment{definition}{\vspace{\parskip}}
\AtBeginEnvironment{example}{\vspace{\parskip}}
\AtBeginEnvironment{lemma}{\vspace{\parskip}}
\AtBeginEnvironment{theorem}{\vspace{\parskip}}
\AtBeginEnvironment{corollary}{\vspace{\parskip}}
\AtBeginEnvironment{proposition}{\vspace{\parskip}}
\AtBeginEnvironment{remark}{\vspace{\parskip}}
\AtBeginEnvironment{exercise}{\vspace{\parskip}}

% Compact section spacing
\usepackage{titlesec}
\titlespacing*{\section}{0pt}{6pt}{3pt}
\titlespacing*{\subsection}{0pt}{4pt}{2pt}
\titlespacing*{\subsubsection}{0pt}{3pt}{1pt}

\begin{document}

% Compact title (no \maketitle top-padding)
\begingroup
\centering
{\LARGE\bfseries Notes on the Modern Theory of Linear Models\par}
\vspace{4pt}
\endgroup

%\tableofcontents
%\newpage

\noindent\textbf{Notation and conventions.}\; Throughout, $\|\bv\| = \|\bv\|_2 = (\sum_i v_i^2)^{1/2}$ is the Euclidean norm, $\opnorm{\bA} = \sup_{\|\bv\|=1}\|\bA\bv\|$ is the operator (spectral) norm, $\fnorm{\bA} = (\tr(\bA'\bA))^{1/2}$ is the Frobenius norm, and $\sigma_{\min}(\bA)$, $\sigma_{\max}(\bA)$ denote the smallest and largest singular values of $\bA$. For a symmetric matrix $\bA$, $\lambda_{\min}(\bA)$ and $\lambda_{\max}(\bA)$ denote its extreme eigenvalues; then $\opnorm{\bA} = \max(|\lambda_{\min}(\bA)|, |\lambda_{\max}(\bA)|)$. We write $a \lesssim b$ to mean $a \leqslant Cb$ for a universal constant $C > 0$. For two symmetric matrices, $\bA \preceq \bA'$ means $\bA' - \bA$ is positive semi-definite. We write $a \wedge b = \min(a,b)$ and $a \vee b = \max(a,b)$. The standard Gaussian random vector in $\Real^n$ is denoted $\mathbf{g} \sim N(\bzero, \bI_n)$. In Part~3, $p$ replaces $k+1$ for the total number of regressors (including the intercept, if present).

\bigskip\noindent\textbf{Prerequisites.}\; This document assumes familiarity with the material in \emph{Notes on the Classical Theory of Linear Models}. The primary references are R.\ Vershynin, \emph{High-Dimensional Probability} (HDP, 2018), and M.\,J.\ Wainwright, \emph{High-Dimensional Statistics} (HDS, 2019).

%=============================================================================
%\newpage
\part{The Toolkit: Sub-Gaussian and Sub-Exponential Variables}\label{part:toolkit}
%=============================================================================

%=============================================================================
\section{Sub-Gaussian Random Variables}\label{sec:subgaussian}
%=============================================================================

\subsection{Motivation: Why Go Beyond Normality?}

\begin{remark}
In the classical linear model (Assumption~(A5)), exact inference rests on the normality of $\be$. The Gauss--Markov result (Theorem~2.6 in the classical notes) needs only (A1)--(A4), but the pivotal distributions $t_{n-k-1}$ and $F_{k,\,n-k-1}$ demand Gaussian errors. The modern programme asks: \emph{what minimal distributional property replaces normality to yield sharp, non-asymptotic concentration?} The answer is the sub-Gaussian condition: the tails of the random variable decay at least as fast as those of a Gaussian. This class includes all bounded random variables, all Gaussian random variables, and many other distributions encountered in practice.
\end{remark}

\subsection{Definition and Equivalent Characterizations}

\begin{definition}[Sub-Gaussian Random Variable]\label{def:subgaussian}
A real-valued random variable $X$ with $\expc[X] = 0$ is called \textbf{sub-Gaussian} if there exists a constant $K > 0$ such that
\begin{equation}\label{eq:subgaussian-mgf}
\expc\!\left[e^{tX}\right] \leqslant e^{K^2 t^2/2} \quad \text{for all } t \in \Real.
\end{equation}
The \textbf{sub-Gaussian norm} (or $\psi_2$-norm) of $X$ is
\begin{equation}\label{eq:psi2}
\|X\|_{\psi_2} = \inf\!\left\{t > 0 : \expc\!\left[e^{X^2/t^2}\right] \leqslant 2\right\}.
\end{equation}
\end{definition}

\begin{theorem}[Equivalent Characterizations of Sub-Gaussianity]\label{thm:subgaussian-equiv}
Let $X$ be a random variable with $\expc[X] = 0$. The following properties are equivalent, in the sense that if one holds with constant $K_i$, then each of the others holds with a constant $K_j = C_{ij} K_i$ for a universal factor $C_{ij} > 0$:
\begin{enumerate}[label=\textup{(\roman*)}]
\item\label{sg:mgf} \textbf{MGF bound}: $\expc[e^{tX}] \leqslant e^{K_1^2 t^2/2}$ for all $t \in \Real$.
\item\label{sg:tail} \textbf{Tail bound}: $\prb(|X| \geqslant t) \leqslant 2\,e^{-t^2/(2K_2^2)}$ for all $t \geqslant 0$.
\item\label{sg:moment} \textbf{Moment bound}: $(\expc[|X|^p])^{1/p} \leqslant K_3\sqrt{p}$ for all $p \geqslant 1$.
\item\label{sg:psi2} \textbf{$\psi_2$-norm finite}: $\|X\|_{\psi_2} \leqslant K_4$.
\end{enumerate}
Moreover, $\|X\|_{\psi_2}$ is equivalent (up to universal constants) to each of the optimal $K_i$.
\end{theorem}

\begin{proof}
We prove the cycle \ref{sg:mgf}~$\Rightarrow$~\ref{sg:tail}~$\Rightarrow$~\ref{sg:moment}~$\Rightarrow$~\ref{sg:psi2}~$\Rightarrow$~\ref{sg:mgf}.

\medskip\noindent\textbf{\ref{sg:mgf}~$\Rightarrow$~\ref{sg:tail}.} By Markov's inequality and~\ref{sg:mgf}, for any $t > 0$ and $\lambda > 0$:
\[
\prb(X \geqslant t) = \prb(e^{\lambda X} \geqslant e^{\lambda t}) \leqslant e^{-\lambda t}\,\expc[e^{\lambda X}] \leqslant e^{-\lambda t + K_1^2\lambda^2/2}.
\]
Optimizing over $\lambda$: set $\frac{d}{d\lambda}(-\lambda t + K_1^2\lambda^2/2) = -t + K_1^2\lambda = 0$, giving $\lambda^* = t/K_1^2$. Substituting:
\[
\prb(X \geqslant t) \leqslant e^{-t^2/(2K_1^2)}.
\]
Applying the same argument to $-X$ (which also satisfies~\ref{sg:mgf} with the same constant, since $\expc[e^{-tX}] = \expc[e^{t(-X)}]$ and $(-X)$ has mean zero):
\[
\prb(|X| \geqslant t) \leqslant 2\,e^{-t^2/(2K_1^2)}.
\]
So $K_2 = K_1$ works.

\medskip\noindent\textbf{\ref{sg:tail}~$\Rightarrow$~\ref{sg:moment}.} For any $p \geqslant 1$, the layer-cake formula gives:
\begin{align}
\expc[|X|^p] &= \int_0^\infty p\,t^{p-1}\prb(|X| \geqslant t)\,\dd t \leqslant \int_0^\infty p\,t^{p-1} \cdot 2\,e^{-t^2/(2K_2^2)}\,\dd t. \label{eq:layercake}
\end{align}
Substituting $s = t^2/(2K_2^2)$, so $t = K_2\sqrt{2s}$ and $\dd t = K_2\,\dd s/\sqrt{2s}$:
\[
\expc[|X|^p] \leqslant 2p \int_0^\infty (K_2\sqrt{2s})^{p-1}\,e^{-s}\,\frac{K_2}{\sqrt{2s}}\,\dd s = p\,(K_2\sqrt{2})^p\,\Gamma(p/2).
\]
For $p \geqslant 2$ we use the bound $\Gamma(\alpha) \leqslant 2\alpha^{\alpha-1}$, valid for $\alpha \geqslant 1$. To see it, write $\Gamma(\alpha) = \int_0^\infty t^{\alpha-1}e^{-t}\,\dd t$ and bound $t^{\alpha-1}e^{-t/2}$ by its maximum over $t \geqslant 0$, attained at $t = 2(\alpha-1)$ with value $(2(\alpha-1)/e)^{\alpha-1}$; then
\[
\Gamma(\alpha) = \int_0^\infty \!\big(t^{\alpha-1}e^{-t/2}\big)e^{-t/2}\,\dd t \leqslant \left(\frac{2(\alpha-1)}{e}\right)^{\!\alpha-1}\!\int_0^\infty e^{-t/2}\,\dd t = 2\left(\frac{2(\alpha-1)}{e}\right)^{\!\alpha-1} \leqslant 2\alpha^{\alpha-1},
\]
the last step using $\alpha - 1 \leqslant \alpha$ and $2/e < 1$. Thus $\Gamma(p/2) \leqslant 2(p/2)^{p/2-1}$, and substituting:
\[
\expc[|X|^p] \leqslant p\,(K_2\sqrt{2})^p\cdot 2(p/2)^{p/2-1} = 4\,(K_2^2)^{p/2}\,p^{p/2}.
\]
Taking $p$-th roots: $(\expc[|X|^p])^{1/p} \leqslant 4^{1/p}K_2\sqrt{p} \leqslant 2K_2\sqrt{p}$ for $p \geqslant 2$. For $p = 1$: $\expc[|X|] \leqslant 2\int_0^\infty e^{-t^2/(2K_2^2)}\,\dd t = K_2\sqrt{2\pi} \leqslant 3K_2$. So $K_3 = 3K_2$ works uniformly for all $p \geqslant 1$.

\medskip\noindent\textbf{\ref{sg:moment}~$\Rightarrow$~\ref{sg:psi2}.} From~\ref{sg:moment}, $\expc[|X|^{2k}] \leqslant (K_3\sqrt{2k})^{2k} = K_3^{2k}(2k)^k$. Using the Taylor expansion:
\[
\expc\!\left[e^{X^2/t^2}\right] = \sum_{k=0}^{\infty} \frac{\expc[|X|^{2k}]}{k!\,t^{2k}} \leqslant 1 + \sum_{k=1}^{\infty} \frac{K_3^{2k}(2k)^k}{k!\,t^{2k}}.
\]
Since $k! \geqslant (k/e)^k$ (Stirling's lower bound), $k^k/k! \leqslant e^k$. Hence $(2k)^k/k! = 2^k\cdot k^k/k! \leqslant (2e)^k$, and:
\[
\frac{K_3^{2k}(2k)^k}{k!\,t^{2k}} \leqslant \left(\frac{2eK_3^2}{t^2}\right)^k.
\]
Choosing $t^2 = 4eK_3^2$: each term is $\leqslant 2^{-k}$, so $\expc[e^{X^2/t^2}] \leqslant 1 + 1 = 2$. Hence $\|X\|_{\psi_2} \leqslant 2\sqrt{e}\,K_3$.

\medskip\noindent\textbf{\ref{sg:psi2}~$\Rightarrow$~\ref{sg:mgf}.} Let $K = \|X\|_{\psi_2}$, so $\expc[e^{X^2/K^2}] \leqslant 2$.

\textbf{Step 1: Obtain a tail bound.} By Markov's inequality: for $u \geqslant 0$,
\[
\prb(|X| \geqslant u) = \prb(e^{X^2/K^2} \geqslant e^{u^2/K^2}) \leqslant e^{-u^2/K^2}\,\expc[e^{X^2/K^2}] \leqslant 2e^{-u^2/K^2}.
\]
This is the tail bound~\ref{sg:tail} with $K_2 = K/\sqrt{2}$.

\textbf{Step 2: Obtain moment bounds.} By the implication \ref{sg:tail}~$\Rightarrow$~\ref{sg:moment} (already proved), $(\expc[|X|^p])^{1/p} \leqslant CK\sqrt{p}$ for all $p \geqslant 1$, hence $|\expc[X^p]| \leqslant \expc[|X|^p] \leqslant (CK)^p p^{p/2}$.

\textbf{Step 3: Bound the MGF via Taylor series.} Since $\expc[X] = 0$:
\[
\expc[e^{tX}] = 1 + \sum_{p=2}^{\infty}\frac{t^p\,\expc[X^p]}{p!} \leqslant 1 + \sum_{p=2}^{\infty}\frac{|t|^p(CK)^p\,p^{p/2}}{p!}.
\]
From Stirling's lower bound $p! \geqslant (p/e)^p$, we obtain $p^{p/2}/p! \leqslant e^p/p^{p/2}$, so
\[
\frac{(CK|t|)^p\,p^{p/2}}{p!} \leqslant \frac{(CK|t|)^p\,e^p}{p^{p/2}} = \left(\frac{CeK|t|}{\sqrt{p}}\right)^p.
\]
For $p \geqslant 2$ and $|t| \leqslant 1/(2CeK)$: $CeK|t|/\sqrt{p} \leqslant CeK|t| \leqslant 1/2$, so each term is $\leqslant 2^{-p}$. Hence:
\[
\expc[e^{tX}] \leqslant 1 + \sum_{p=2}^\infty 2^{-p} = 1 + 1/2 < 2 \leqslant e^1, \quad \text{for } |t| \leqslant \frac{1}{2CeK}.
\]

\textbf{Step 4: Extend to all $t$.} For $|t| > 1/(2CeK)$, we use the tail bound from Step~1. Split the MGF:
\[
\expc[e^{tX}] = \expc[e^{tX}\mathbf{1}_{|X| \leqslant R}] + \expc[e^{tX}\mathbf{1}_{|X| > R}]
\]
with $R = K^2|t|$. For the first term: $e^{tX}\mathbf{1}_{|X| \leqslant R} \leqslant e^{|t|R} = e^{K^2t^2}$. For the second:
\begin{align*}
\expc[e^{tX}\mathbf{1}_{|X|>R}] &\leqslant \expc[e^{|t||X|}\mathbf{1}_{|X|>R}] = |t|\int_R^\infty e^{|t|u}\prb(|X| \geqslant u)\,\dd u + e^{|t|R}\prb(|X| \geqslant R)\\
&\leqslant |t|\int_R^\infty e^{|t|u}\cdot 2e^{-u^2/K^2}\,\dd u + 2e^{|t|R - R^2/K^2}.
\end{align*}
In the exponent: $|t|u - u^2/K^2 = -(u/K - K|t|/2)^2 + K^2t^2/4$. So all terms are bounded by $Ce^{K^2t^2/4}\cdot(\text{polynomial in }|t|K)$, giving $\expc[e^{tX}] \leqslant e^{C'K^2t^2}$ for a universal $C'$.

Combining Steps 3--4: $\expc[e^{tX}] \leqslant e^{C''K^2t^2/2}$ for all $t \in \Real$, with $K_1 = \sqrt{2C''}K_4$.
\end{proof}

\begin{example}[Sub-Gaussian Random Variables]\label{ex:subgaussian}
\leavevmode
\begin{enumerate}[label=\textup{(\arabic*)}]
\item \textbf{Gaussian.} If $X \sim N(0, \sigma^2)$, then $\expc[e^{tX}] = e^{\sigma^2 t^2/2}$, so $X$ is sub-Gaussian with $K_1 = \sigma$. Moreover, $\|X\|_{\psi_2} = C\sigma$ for a universal $C$.

\item \textbf{Bounded (Rademacher).} Let $X$ take values $\pm 1$ with equal probability. Then $\expc[e^{tX}] = \cosh(t) \leqslant e^{t^2/2}$, so $X$ is sub-Gaussian with $K_1 = 1$.

\item \textbf{Bounded (Hoeffding's lemma).} If $a \leqslant X \leqslant b$ and $\expc[X] = 0$, then $\expc[e^{tX}] \leqslant e^{(b-a)^2t^2/8}$, so $X$ is sub-Gaussian with $K_1 = (b-a)/2$. The proof is given in Lemma~\ref{lem:hoeffding-lemma}.

\item \textbf{Non-example (Exponential).} If $X \sim \mathrm{Exp}(1) - 1$ (centered), then $\expc[e^{tX}] = e^{-t}/(1-t)$ for $t < 1$, which grows faster than any $e^{Kt^2}$ as $t \to 1^-$. Hence $X$ is \emph{not} sub-Gaussian; it is sub-exponential (Definition~\ref{def:subexponential}).
\end{enumerate}
\end{example}

\subsection{Properties of Sub-Gaussian Random Variables}

\begin{proposition}[Closure Properties]\label{prop:sg-closure}
\leavevmode
\begin{enumerate}[label=\textup{(\roman*)}]
\item If $X$ is sub-Gaussian and $a \in \Real$, then $aX$ is sub-Gaussian with $\|aX\|_{\psi_2} = |a|\,\|X\|_{\psi_2}$.
\item If $X_1, \ldots, X_n$ are independent sub-Gaussian random variables with $\expc[X_i] = 0$, then $S = \sum_{i=1}^n X_i$ is sub-Gaussian with $\|S\|_{\psi_2}^2 \leqslant C\sum_{i=1}^n \|X_i\|_{\psi_2}^2$.
\end{enumerate}
\end{proposition}

\begin{proof}
\textbf{(i).} Immediate from $\expc[e^{(aX)^2/t^2}] = \expc[e^{X^2/(t/|a|)^2}]$, which is $\leqslant 2$ when $t/|a| \geqslant \|X\|_{\psi_2}$, i.e., $t \geqslant |a|\,\|X\|_{\psi_2}$.

\textbf{(ii).} By independence:
\[
\expc\!\left[e^{tS}\right] = \prod_{i=1}^n \expc\!\left[e^{tX_i}\right] \leqslant \prod_{i=1}^n e^{K_i^2 t^2/2} = \exp\!\left(\frac{t^2}{2}\sum_{i=1}^n K_i^2\right),
\]
where $K_i = C\|X_i\|_{\psi_2}$ from the MGF characterization. Hence $S$ is sub-Gaussian with $K^2 = \sum K_i^2$, and $\|S\|_{\psi_2}^2 \leqslant C'\sum\|X_i\|_{\psi_2}^2$.
\end{proof}

\subsection{Variance of Sub-Gaussian Variables}

\begin{proposition}\label{prop:sg-variance}
If $X$ is sub-Gaussian with $\expc[X] = 0$, then $\var(X) = \expc[X^2] \leqslant C\|X\|_{\psi_2}^2$ for a universal constant $C$.
\end{proposition}

\begin{proof}
From the moment bound (Theorem~\ref{thm:subgaussian-equiv}\ref{sg:moment}) with $p = 2$: $(\expc[X^2])^{1/2} \leqslant K_3\sqrt{2}$. Since $K_3 \leqslant C\|X\|_{\psi_2}$, squaring gives $\expc[X^2] \leqslant 2C^2\|X\|_{\psi_2}^2$.
\end{proof}

%=============================================================================
\section{Hoeffding's Inequality}\label{sec:hoeffding}
%=============================================================================

\subsection{Hoeffding's Lemma}

\begin{lemma}[Hoeffding's Lemma]\label{lem:hoeffding-lemma}
Let $X$ be a random variable with $\expc[X] = 0$ and $a \leqslant X \leqslant b$ almost surely. Then for all $t \in \Real$:
\begin{equation}\label{eq:hoeffding-lemma}
\expc[e^{tX}] \leqslant \exp\!\left(\frac{t^2(b-a)^2}{8}\right).
\end{equation}
\end{lemma}

\begin{proof}
If $a = b$, then $\expc[X] = 0$ forces $X = a = b = 0$ almost surely, and both sides of~\eqref{eq:hoeffding-lemma} equal $1$; assume henceforth $a < b$.

\textbf{Step 1: Convexity bound.} Since $e^{tx}$ is convex in $x$, for any $x \in [a, b]$:
\[
e^{tx} \leqslant \frac{b - x}{b - a}\,e^{ta} + \frac{x - a}{b - a}\,e^{tb}.
\]
Taking expectations and using $\expc[X] = 0$:
\begin{equation}\label{eq:hoeffding-step1}
\expc[e^{tX}] \leqslant \frac{b}{b - a}\,e^{ta} + \frac{-a}{b - a}\,e^{tb} = \frac{b\,e^{ta} - a\,e^{tb}}{b - a}.
\end{equation}

\textbf{Step 2: Logarithmic bound.} Let $h = t(b-a)$, $\theta = -a/(b-a) \in [0,1]$. Then $1 - \theta = b/(b-a)$. Rewrite~\eqref{eq:hoeffding-step1}:
\begin{align*}
\expc[e^{tX}] &\leqslant (1-\theta)\,e^{-\theta h} + \theta\,e^{(1-\theta)h} = e^{-\theta h}\bigl[(1-\theta) + \theta\,e^{h}\bigr].
\end{align*}
Define $\varphi(h) = -\theta h + \ln\bigl[(1-\theta) + \theta\,e^h\bigr]$. Then $\expc[e^{tX}] \leqslant e^{\varphi(h)}$.

\textbf{Step 3: Taylor expansion of $\varphi$.} We have $\varphi(0) = 0$ and $\varphi'(0) = -\theta + \theta = 0$. Computing the second derivative:
\[
\varphi''(h) = \frac{\theta(1-\theta)\,e^h}{[(1-\theta) + \theta\,e^h]^2}.
\]
Let $u = \theta\,e^h/[(1-\theta) + \theta\,e^h] \in (0,1)$. Then $\varphi''(h) = u(1-u) \leqslant 1/4$ (since $u(1-u) \leqslant 1/4$ for $u \in [0,1]$). By Taylor's theorem with remainder:
\[
\varphi(h) = \varphi(0) + h\varphi'(0) + \frac{h^2}{2}\varphi''(\xi) \leqslant \frac{h^2}{8}
\]
for some $\xi$ between $0$ and $h$. Hence $\expc[e^{tX}] \leqslant e^{h^2/8} = e^{t^2(b-a)^2/8}$.
\end{proof}

\subsection{Hoeffding's Inequality}

\begin{theorem}[Hoeffding's Inequality]\label{thm:hoeffding}
Let $X_1, \ldots, X_n$ be independent random variables with $\expc[X_i] = 0$ and $a_i \leqslant X_i \leqslant b_i$ almost surely for each $i$. Then for all $t \geqslant 0$:
\begin{equation}\label{eq:hoeffding}
\boxed{\prb\!\left(\left|\sum_{i=1}^n X_i\right| \geqslant t\right) \leqslant 2\exp\!\left(-\frac{2t^2}{\sum_{i=1}^n (b_i - a_i)^2}\right).}
\end{equation}
\end{theorem}

\begin{proof}
\textbf{Step 1: One-sided bound via Chernoff.} For $\lambda > 0$:
\begin{align}
\prb\!\left(\sum_{i=1}^n X_i \geqslant t\right) &= \prb\!\left(e^{\lambda\sum X_i} \geqslant e^{\lambda t}\right) \leqslant e^{-\lambda t}\,\expc\!\left[e^{\lambda\sum X_i}\right] \notag\\
&= e^{-\lambda t}\prod_{i=1}^n \expc[e^{\lambda X_i}] \leqslant e^{-\lambda t}\prod_{i=1}^n e^{\lambda^2(b_i - a_i)^2/8}, \label{eq:hoeffding-chernoff}
\end{align}
where the last inequality uses Hoeffding's Lemma (Lemma~\ref{lem:hoeffding-lemma}) applied to each $X_i$.

\textbf{Step 2: Optimize $\lambda$.} From~\eqref{eq:hoeffding-chernoff}:
\[
\prb\!\left(\sum X_i \geqslant t\right) \leqslant \exp\!\left(-\lambda t + \frac{\lambda^2}{8}\sum_{i=1}^n(b_i - a_i)^2\right).
\]
Let $\sigma_H^2 = \frac{1}{4}\sum(b_i - a_i)^2$. Setting $\frac{d}{d\lambda}(-\lambda t + \lambda^2\sigma_H^2/2) = -t + \lambda\sigma_H^2 = 0$ gives $\lambda^* = t/\sigma_H^2$:
\[
\prb\!\left(\sum X_i \geqslant t\right) \leqslant \exp\!\left(-\frac{t^2}{2\sigma_H^2}\right) = \exp\!\left(-\frac{2t^2}{\sum(b_i-a_i)^2}\right).
\]

\textbf{Step 3: Two-sided bound.} Apply the same argument to $-X_1, \ldots, -X_n$ to get $\prb(\sum X_i \leqslant -t) \leqslant \exp(-2t^2/\sum(b_i-a_i)^2)$. By the union bound, $\prb(|\sum X_i| \geqslant t) \leqslant 2\exp(-2t^2/\sum(b_i-a_i)^2)$.
\end{proof}

\begin{corollary}[Hoeffding for i.i.d.\ Bounded Variables]\label{cor:hoeffding-iid}
Let $X_1, \ldots, X_n$ be i.i.d.\ with $\expc[X_i] = \mu$ and $a \leqslant X_i \leqslant b$. Then for $\bar{X} = \frac{1}{n}\sum X_i$:
\begin{equation}\label{eq:hoeffding-mean}
\prb\!\left(|\bar{X} - \mu| \geqslant t\right) \leqslant 2\exp\!\left(-\frac{2nt^2}{(b-a)^2}\right).
\end{equation}
\end{corollary}

\begin{proof}
Apply Theorem~\ref{thm:hoeffding} to $Y_i = X_i - \mu$ (so $\expc[Y_i] = 0$, $a - \mu \leqslant Y_i \leqslant b - \mu$, and $b_i - a_i = b - a$ for all $i$). Then $\sum Y_i = n(\bar{X} - \mu)$, and the bound gives $\prb(|n(\bar{X}-\mu)| \geqslant nt) \leqslant 2\exp(-2n^2t^2/(n(b-a)^2))$.
\end{proof}

%=============================================================================
\section{Sub-Gaussian Concentration for Sums and Quadratic Forms}\label{sec:sg-conc}
%=============================================================================

\subsection{General Sub-Gaussian Hoeffding Inequality}

\begin{theorem}[General Hoeffding Inequality for Sub-Gaussians]\label{thm:hoeffding-sg}
Let $X_1, \ldots, X_n$ be independent, mean-zero, sub-Gaussian random variables. Then for all $t \geqslant 0$:
\begin{equation}\label{eq:hoeffding-sg}
\prb\!\left(\left|\sum_{i=1}^n X_i\right| \geqslant t\right) \leqslant 2\exp\!\left(-\frac{ct^2}{\sum_{i=1}^n \|X_i\|_{\psi_2}^2}\right),
\end{equation}
where $c > 0$ is a universal constant.
\end{theorem}

\begin{proof}
By the MGF characterization (Theorem~\ref{thm:subgaussian-equiv}\ref{sg:mgf}), each $X_i$ satisfies $\expc[e^{\lambda X_i}] \leqslant e^{C^2\|X_i\|_{\psi_2}^2\lambda^2/2}$. The proof then follows the Chernoff--optimize pattern of Theorem~\ref{thm:hoeffding}, with $(b_i - a_i)^2/4$ replaced by $C^2\|X_i\|_{\psi_2}^2$:
\begin{align*}
\prb\!\left(\sum X_i \geqslant t\right) &\leqslant \exp\!\left(-\lambda t + \frac{C^2\lambda^2}{2}\sum \|X_i\|_{\psi_2}^2\right).
\end{align*}
Setting $\lambda^* = t/(C^2\sum\|X_i\|_{\psi_2}^2)$ gives $\prb(\sum X_i \geqslant t) \leqslant \exp(-t^2/(2C^2\sum\|X_i\|_{\psi_2}^2))$, and the two-sided bound follows by the union bound with $c = 1/(2C^2)$.
\end{proof}

\subsection{Sub-Gaussian Random Vectors}

\begin{definition}[Sub-Gaussian Random Vector]\label{def:sg-vector}
A random vector $\bx \in \Real^p$ is \textbf{sub-Gaussian} if all one-dimensional marginals are sub-Gaussian. Formally, there exists $K > 0$ such that
\begin{equation}\label{eq:sg-vector}
\|\ip{\bx}{\bv}\|_{\psi_2} \leqslant K\,\|\bv\| \quad \text{for all } \bv \in \Real^p.
\end{equation}
The smallest such $K$ is the \textbf{sub-Gaussian norm} of $\bx$, denoted $\|\bx\|_{\psi_2}$.
\end{definition}

\begin{example}\label{ex:sg-vector}
\leavevmode
\begin{enumerate}[label=\textup{(\arabic*)}]
\item If $\bx \sim N(\bzero, \bSigma)$, then $\ip{\bx}{\bv} \sim N(0, \bv'\bSigma\bv)$ has $\|\ip{\bx}{\bv}\|_{\psi_2} = C(\bv'\bSigma\bv)^{1/2} \leqslant C\|\bSigma\|_{\mathrm{op}}^{1/2}\|\bv\|$. So $\bx$ is sub-Gaussian with $\|\bx\|_{\psi_2} \leqslant C\opnorm{\bSigma}^{1/2}$.
\item If $\bx = (\varepsilon_1, \ldots, \varepsilon_p)'$ where $\varepsilon_i$ are independent Rademacher, then $\ip{\bx}{\bv} = \sum v_i\varepsilon_i$ has $\|\ip{\bx}{\bv}\|_{\psi_2} \leqslant C\|\bv\|$ by Proposition~\ref{prop:sg-closure}(ii).
\end{enumerate}
\end{example}

\subsection{Hanson--Wright Inequality for Quadratic Forms}

In this subsection we prove the Hanson--Wright concentration inequality for quadratic forms $\bx'\bA\bx$, where $\bx$ has independent sub-Gaussian coordinates. The result drives nearly every application in Parts~3--4. We give a complete, self-contained proof, with no step deferred to an external reference.

The presentation has three parts. (i)~The Gaussian case (Theorem~\ref{thm:hanson-wright-gaussian}) is proved by computing the moment generating function (MGF) of a weighted sum of independent $\chi^2_1$ variables, using the spectral decomposition of $\bA$. (ii)~The sub-Gaussian case (Theorem~\ref{thm:hanson-wright}) is proved in eight steps. The two substantive ingredients beyond the Gaussian argument --- a \emph{decoupling lemma} and an \emph{MGF lemma} for sub-Gaussian quadratic forms --- are stated and proved in full; the MGF lemma uses the Hubbard--Stratonovich identity to reduce the sub-Gaussian quadratic-form MGF directly to the Gaussian one. (iii)~We apply the result.

\subsubsection*{Gaussian case (full proof)}

\begin{theorem}[Hanson--Wright, Gaussian case]\label{thm:hanson-wright-gaussian}
Let $\bg \sim N(\bzero, \bI_n)$ and let $\bA$ be any $n \times n$ matrix. Then for all $t \geqslant 0$:
\begin{equation}\label{eq:hanson-wright-gaussian}
\prb\!\left(\left|\bg'\bA\bg - \tr(\bA)\right| \geqslant t\right) \leqslant 2\exp\!\left[-c\min\!\left(\frac{t^2}{\fnorm{\bA}^2},\; \frac{t}{\opnorm{\bA}}\right)\right],
\end{equation}
where $c = 1/8$ suffices.
\end{theorem}

\begin{proof}
Since $\bg'\bA\bg = \bg'\bA_s\bg$ and $\tr(\bA) = \tr(\bA_s)$ for the symmetric part $\bA_s = (\bA + \bA')/2$, while $\fnorm{\bA_s} \leqslant \fnorm{\bA}$ and $\opnorm{\bA_s} \leqslant \opnorm{\bA}$ (as $\fnorm{\bA'} = \fnorm{\bA}$, $\opnorm{\bA'} = \opnorm{\bA}$) make the right-hand side of~\eqref{eq:hanson-wright-gaussian} only larger, it suffices to prove the bound for $\bA_s$; we may thus assume $\bA$ is symmetric. Let $\bA = \bU\bLambda\bU'$ with $\bLambda = \diag(\lambda_1, \ldots, \lambda_n)$. Set $\widetilde\bg = \bU'\bg$, which is again $N(\bzero, \bI_n)$ since $\bU$ is orthogonal. Then
\[
\bg'\bA\bg = \widetilde\bg'\bLambda\widetilde\bg = \sum_{k=1}^n \lambda_k \widetilde g_k^2,
\]
a weighted sum of independent $\chi^2_1$ variables.

\textbf{MGF.} For any $\mu$ with $\mu\lambda_k < 1/2$ for all $k$,
\[
\expc[e^{\mu \widetilde g_k^2}] = \int_{-\infty}^\infty \frac{e^{\mu u^2 - u^2/2}}{\sqrt{2\pi}}\,\dd u = \frac{1}{\sqrt{1 - 2\mu}}.
\]
Hence, for $|\mu| \opnorm{\bA} < 1/2$,
\begin{equation}\label{eq:hw-gauss-mgf}
\expc[e^{\mu (\bg'\bA\bg - \tr(\bA))}] = \prod_{k=1}^n \frac{e^{-\mu \lambda_k}}{\sqrt{1 - 2\mu\lambda_k}}.
\end{equation}

\textbf{Bound the MGF.} For $|u| \leqslant 1/4$, the elementary inequality $-\ln(1 - 2u) - 2u \leqslant 4u^2$ holds (Taylor expansion: $-\ln(1-2u) - 2u = 2u^2 + \tfrac{8}{3}u^3 + 4u^4 + \cdots \leqslant 2u^2/(1-2|u|) \leqslant 4u^2$ for $|u| \leqslant 1/4$). Applying this with $u = \mu\lambda_k$ for all $k$ (which requires $\mu\opnorm{\bA} \leqslant 1/4$) and summing the halved version $-\tfrac{1}{2}\ln(1 - 2u) - u \leqslant 2u^2$:
\[
\sum_{k=1}^n \left[-\tfrac{1}{2}\ln(1 - 2\mu\lambda_k) - \mu\lambda_k\right] \leqslant 2\mu^2 \sum_{k=1}^n \lambda_k^2 = 2\mu^2 \fnorm{\bA}^2,
\]
so
\begin{equation}\label{eq:hw-gauss-mgf-bound}
\expc[e^{\mu (\bg'\bA\bg - \tr(\bA))}] \leqslant e^{2\mu^2 \fnorm{\bA}^2} \quad \text{for } |\mu| \leqslant \frac{1}{4\opnorm{\bA}}.
\end{equation}

\textbf{Chernoff with constrained optimization.} For $t > 0$, Markov's inequality gives
\[
\prb(\bg'\bA\bg - \tr(\bA) \geqslant t) \leqslant \inf_{0 < \mu \leqslant 1/(4\opnorm{\bA})} \exp(-\mu t + 2\mu^2 \fnorm{\bA}^2).
\]
The unconstrained optimum is $\mu^* = t/(4\fnorm{\bA}^2)$.

\emph{Case 1 (sub-Gaussian regime):} $\mu^* \leqslant 1/(4\opnorm{\bA})$, i.e., $t \leqslant \fnorm{\bA}^2/\opnorm{\bA}$. Substituting $\mu^*$:
\[
\prb(\bg'\bA\bg - \tr(\bA) \geqslant t) \leqslant \exp\!\left(-\frac{t^2}{8\fnorm{\bA}^2}\right).
\]

\emph{Case 2 (sub-exponential regime):} $t > \fnorm{\bA}^2/\opnorm{\bA}$. Set $\mu = 1/(4\opnorm{\bA})$:
\[
-\mu t + 2\mu^2 \fnorm{\bA}^2 = -\frac{t}{4\opnorm{\bA}} + \frac{\fnorm{\bA}^2}{8\opnorm{\bA}^2} \leqslant -\frac{t}{4\opnorm{\bA}} + \frac{t}{8\opnorm{\bA}} = -\frac{t}{8\opnorm{\bA}},
\]
using $\fnorm{\bA}^2/\opnorm{\bA}^2 \leqslant t/\opnorm{\bA}$ in Case~2. Hence
\[
\prb(\bg'\bA\bg - \tr(\bA) \geqslant t) \leqslant \exp\!\left(-\frac{t}{8\opnorm{\bA}}\right).
\]
Combining the two cases gives the one-sided bound with $c = 1/8$. Applying the same argument to $-\bA$ and a union bound yields~\eqref{eq:hanson-wright-gaussian}.
\end{proof}

\subsubsection*{Sub-Gaussian case}

\begin{theorem}[Hanson--Wright Inequality]\label{thm:hanson-wright}
Let $\bx = (X_1, \ldots, X_n)' \in \Real^n$ be a random vector with independent, mean-zero, sub-Gaussian components satisfying $\|X_i\|_{\psi_2} \leqslant K$. Let $\bA$ be an $n \times n$ matrix. Then for all $t \geqslant 0$:
\begin{equation}\label{eq:hanson-wright}
\boxed{\prb\!\left(\left|\bx'\bA\bx - \expc[\bx'\bA\bx]\right| \geqslant t\right) \leqslant 2\exp\!\left[-c\min\!\left(\frac{t^2}{K^4\fnorm{\bA}^2},\; \frac{t}{K^2\opnorm{\bA}}\right)\right],}
\end{equation}
where $c > 0$ is a universal constant.
\end{theorem}

The proof follows the skeleton of the Gaussian case, with two substantive lemmas inserted: a \emph{decoupling lemma} (Step~2) and an \emph{MGF lemma} for sub-Gaussian quadratic forms (Step~4). Both are stated and proved in full. Nothing else is used beyond the Chernoff method, Bernstein's inequality (Theorem~\ref{thm:bernstein}), the sub-Gaussian MGF characterization (Theorem~\ref{thm:subgaussian-equiv}\ref{sg:mgf}), and the Gaussian case just proved (Theorem~\ref{thm:hanson-wright-gaussian}). The argument has eight steps.

\begin{proof}
\textbf{Step 1: Reduce to the zero-diagonal case.} Write $\bA = \bD + \bA_0$, where $\bD$ is the diagonal matrix $\diag(a_{11}, \ldots, a_{nn})$ and $\bA_0$ has zero diagonal. Then
\[
\bx'\bA\bx - \expc[\bx'\bA\bx] = \underbrace{\sum_{i=1}^n a_{ii}(X_i^2 - \expc[X_i^2])}_{\text{diagonal part}} + \underbrace{\bx'\bA_0\bx}_{\text{off-diagonal part}},
\]
since $\expc[\bx'\bA_0\bx] = \sum_{i \neq j}(\bA_0)_{ij}\expc[X_i]\expc[X_j] = 0$ by independence and mean zero.

\emph{Diagonal part.} The variables $Y_i = X_i^2 - \expc[X_i^2]$ are independent and mean-zero. By Proposition~\ref{prop:sg-square}, $X_i^2$ is sub-exponential with $\|X_i^2\|_{\psi_1} \leqslant C_0 K^2$; subtracting the constant $\expc[X_i^2]$ changes the $\psi_1$-norm by at most a constant factor, since $|\expc[X_i^2]| \leqslant C_0 K^2$ by Proposition~\ref{prop:sg-variance}. Hence $\|Y_i\|_{\psi_1} \leqslant C_0 K^2$. The summands $a_{ii}Y_i$ then satisfy $\sum_i \|a_{ii}Y_i\|_{\psi_1}^2 \leqslant C_0^2 K^4\fnorm{\bD}^2$ and $\max_i \|a_{ii}Y_i\|_{\psi_1} \leqslant C_0 K^2\opnorm{\bD}$, so Bernstein's inequality (Theorem~\ref{thm:bernstein}) gives
\begin{equation}\label{eq:hw-diagonal}
\prb\!\left(\left|\sum_{i=1}^n a_{ii}Y_i\right| \geqslant t\right) \leqslant 2\exp\!\left[-c_1\min\!\left(\frac{t^2}{K^4\fnorm{\bD}^2},\;\frac{t}{K^2\opnorm{\bD}}\right)\right],
\end{equation}
with $\fnorm{\bD} \leqslant \fnorm{\bA}$ and $\opnorm{\bD} = \max_i|a_{ii}| \leqslant \opnorm{\bA}$. It remains to bound the off-diagonal part $\bx'\bA_0\bx$, which occupies Steps~2--7.

\textbf{Step 2: Decoupling lemma.} For a subset $I \subseteq \{1, \ldots, n\}$, write $I^c$ for its complement and let $\bA_0^{(I)}$ be the \emph{masked matrix} with entries $(\bA_0^{(I)})_{ij} = (\bA_0)_{ij}$ when $i \in I$ and $j \in I^c$, and $(\bA_0^{(I)})_{ij} = 0$ otherwise. Let $\by$ be an independent copy of $\bx$.

\begin{quote}
\emph{Decoupling lemma. For every convex function $F : \Real \to \Real$,}
\begin{equation}\label{eq:hw-decoupling}
\expc[F(\bx'\bA_0\bx)] \leqslant \frac{1}{2^n}\sum_{I \subseteq \{1,\ldots,n\}} \expc\!\left[F\!\left(4\,\bx'\bA_0^{(I)}\by\right)\right].
\end{equation}
\emph{Moreover, $\fnorm{\bA_0^{(I)}} \leqslant \fnorm{\bA_0}$ and $\opnorm{\bA_0^{(I)}} \leqslant \opnorm{\bA_0}$ for every $I$.}
\end{quote}

\emph{Proof.} Let $\delta_1, \ldots, \delta_n$ be independent $\mathrm{Bernoulli}(1/2)$ variables, independent of $\bx$ and $\by$, and set $I = \{i : \delta_i = 1\}$. For any ordered pair $i \neq j$, independence gives $\prb(i \in I,\, j \in I^c) = \prb(\delta_i = 1)\prb(\delta_j = 0) = 1/4$. Since $\bA_0$ has zero diagonal, $\bx'\bA_0\bx = \sum_{i \neq j}(\bA_0)_{ij}X_iX_j$, so, writing $\expc_I$ for expectation over the $\delta_i$ only,
\[
\expc_I\!\left[\,\sum_{i \in I,\, j \in I^c}(\bA_0)_{ij}X_iX_j\,\right] = \sum_{i \neq j}(\bA_0)_{ij}X_iX_j \cdot \prb(i \in I,\, j \in I^c) = \tfrac{1}{4}\,\bx'\bA_0\bx.
\]
Thus $\bx'\bA_0\bx = \expc_I[\,4\sum_{i \in I, j \in I^c}(\bA_0)_{ij}X_iX_j\,]$. Applying Jensen's inequality to the convex $F$ (conditionally on $\bx$),
\[
F(\bx'\bA_0\bx) \leqslant \expc_I\!\left[F\!\left(4\sum_{i \in I,\, j \in I^c}(\bA_0)_{ij}X_iX_j\right)\right].
\]
Take $\expc_\bx$ on both sides; since $\expc_I$ is an average over the $2^n$ subsets, and the convex $F$ is bounded below by an affine function (so each term is well-defined in $(-\infty, +\infty]$), linearity of expectation over this finite sum gives
\[
\expc_\bx[F(\bx'\bA_0\bx)] \leqslant \frac{1}{2^n}\sum_I \expc_\bx\!\left[F\!\left(4\sum_{i \in I,\, j \in I^c}(\bA_0)_{ij}X_iX_j\right)\right].
\]
Now fix a subset $I$. The sum $\sum_{i \in I, j \in I^c}(\bA_0)_{ij}X_iX_j$ depends on $\bx$ only through the two disjoint --- hence independent --- coordinate blocks $(X_i)_{i \in I}$ and $(X_j)_{j \in I^c}$. Because $\by$ is an independent copy of $\bx$, the pair of blocks $\big((X_i)_{i \in I},\, (X_j)_{j \in I^c}\big)$ has the same joint distribution as $\big((X_i)_{i \in I},\, (Y_j)_{j \in I^c}\big)$. Replacing $X_j$ by $Y_j$ for $j \in I^c$ therefore leaves the distribution of the bilinear sum unchanged:
\[
\sum_{i \in I,\, j \in I^c}(\bA_0)_{ij}X_iX_j \;\overset{d}{=}\; \sum_{i \in I,\, j \in I^c}(\bA_0)_{ij}X_iY_j = \bx'\bA_0^{(I)}\by,
\]
and substituting into the previous display yields~\eqref{eq:hw-decoupling}.

For the norm bounds: $\bA_0^{(I)}$ is obtained from $\bA_0$ by setting some entries to zero, so $\fnorm{\bA_0^{(I)}}^2 = \sum_{i \in I, j \in I^c}(\bA_0)_{ij}^2 \leqslant \fnorm{\bA_0}^2$. Writing $\bP_I$ for the diagonal $0/1$ projection onto the coordinates indexed by $I$ and $\bP_{I^c} = \bI_n - \bP_I$, one has $\bA_0^{(I)} = \bP_I\,\bA_0\,\bP_{I^c}$; since $\opnorm{\bP_I} \leqslant 1$ and $\opnorm{\bP_{I^c}} \leqslant 1$, it follows that $\opnorm{\bA_0^{(I)}} \leqslant \opnorm{\bA_0}$. \hfill$\square$

\textbf{Step 3: Conditional sub-Gaussian MGF of the bilinear form.} Fix a subset $I$ and condition on $\by$. Then $\bx'\bA_0^{(I)}\by = \sum_{i=1}^n X_i\,(\bA_0^{(I)}\by)_i$ is, in the randomness of $\bx$, a sum of independent mean-zero random variables. Each summand is sub-Gaussian with $\|X_i(\bA_0^{(I)}\by)_i\|_{\psi_2} = |(\bA_0^{(I)}\by)_i|\,\|X_i\|_{\psi_2} \leqslant K\,|(\bA_0^{(I)}\by)_i|$. By the sub-Gaussian MGF characterization (Theorem~\ref{thm:subgaussian-equiv}\ref{sg:mgf}), there is a universal constant $C_1 > 0$ with $\expc[e^{s X_i(\bA_0^{(I)}\by)_i} \mid \by] \leqslant \exp(C_1 K^2 (\bA_0^{(I)}\by)_i^2\, s^2/2)$ for every $s \in \Real$. Multiplying over the independent coordinates,
\begin{equation}\label{eq:hw-cond-mgf}
\expc\!\left[e^{\lambda\,\bx'\bA_0^{(I)}\by} \,\middle|\, \by\right] = \prod_{i=1}^n \expc\!\left[e^{\lambda X_i(\bA_0^{(I)}\by)_i} \,\middle|\, \by\right] \leqslant \exp\!\left(\frac{C_1 K^2 \lambda^2 \,\|\bA_0^{(I)}\by\|^2}{2}\right), \quad \lambda \in \Real.
\end{equation}

\textbf{Step 4: MGF lemma for sub-Gaussian quadratic forms.} To take the expectation of~\eqref{eq:hw-cond-mgf} over $\by$ we need an MGF bound for the quadratic form $\|\bA_0^{(I)}\by\|^2 = \by'(\bA_0^{(I)})'\bA_0^{(I)}\by$. We prove the following.

\begin{quote}
\emph{MGF lemma. Let $\by \in \Real^n$ have independent, mean-zero, sub-Gaussian coordinates with $\|Y_i\|_{\psi_2} \leqslant K$, and let $\bB$ be a symmetric positive semidefinite $n \times n$ matrix. Then there are universal constants $C, c_2 > 0$ such that}
\begin{equation}\label{eq:hw-mgf-quadratic}
\expc\!\left[e^{\mu\,\by'\bB\by}\right] \leqslant \exp\!\left(C K^2\mu\,\tr(\bB) + C^2 K^4\mu^2\,\fnorm{\bB}^2\right) \quad \text{for } 0 \leqslant \mu \leqslant \frac{c_2}{K^2\opnorm{\bB}}.
\end{equation}
\end{quote}

\emph{Proof.} The argument has three parts.

\emph{(1) Hubbard--Stratonovich identity.} For any $\bv \in \Real^n$ and $\bg \sim N(\bzero, \bI_n)$,
\[
\expc_\bg\!\left[e^{\ip{\bv}{\bg}}\right] = \prod_{i=1}^n \expc[e^{v_i g_i}] = \prod_{i=1}^n e^{v_i^2/2} = e^{\|\bv\|^2/2}.
\]
Let $\bB^{1/2}$ be the symmetric positive semidefinite square root of $\bB$. Applying the identity with $\bv = \sqrt{2\mu}\,\bB^{1/2}\by$ (here $\mu \geqslant 0$), and noting $\|\bv\|^2/2 = \mu\,\by'\bB^{1/2}\bB^{1/2}\by = \mu\,\by'\bB\by$,
\begin{equation}\label{eq:hw-hubbard}
e^{\mu\,\by'\bB\by} = \expc_\bg\!\left[e^{\sqrt{2\mu}\,\bg'\bB^{1/2}\by}\right].
\end{equation}

\emph{(2) Swap the expectations and use the sub-Gaussian MGF of a linear form.} Both sides of~\eqref{eq:hw-hubbard} are nonnegative, so taking $\expc_\by$ and applying Tonelli's theorem to exchange the order of integration,
\[
\expc_\by\!\left[e^{\mu\,\by'\bB\by}\right] = \expc_\bg\,\expc_\by\!\left[e^{\sqrt{2\mu}\,(\bB^{1/2}\bg)'\by}\right].
\]
Fix $\bg$ and set $\bw = \bB^{1/2}\bg$. The coordinates $Y_i$ are independent, mean-zero, with $\|Y_i\|_{\psi_2} \leqslant K$, so by Theorem~\ref{thm:subgaussian-equiv}\ref{sg:mgf}, $\expc[e^{s Y_i}] \leqslant e^{C_1 K^2 s^2/2}$ with the same universal $C_1$ as in Step~3. Hence
\[
\expc_\by\!\left[e^{\sqrt{2\mu}\,\bw'\by}\right] = \prod_{i=1}^n \expc\!\left[e^{\sqrt{2\mu}\,w_i Y_i}\right] \leqslant \prod_{i=1}^n e^{C_1 K^2 \mu\, w_i^2} = e^{C_1 K^2 \mu\,\|\bw\|^2} = e^{C_1 K^2 \mu\,\bg'\bB\bg},
\]
using $\|\bw\|^2 = \bg'\bB^{1/2}\bB^{1/2}\bg = \bg'\bB\bg$. Therefore
\begin{equation}\label{eq:hw-mgf-reduced}
\expc_\by\!\left[e^{\mu\,\by'\bB\by}\right] \leqslant \expc_\bg\!\left[e^{C_1 K^2 \mu\,\bg'\bB\bg}\right].
\end{equation}

\emph{(3) Apply the Gaussian quadratic-form MGF.} The right-hand side of~\eqref{eq:hw-mgf-reduced} is a Gaussian quadratic-form MGF, controlled by inequality~\eqref{eq:hw-gauss-mgf-bound} of Theorem~\ref{thm:hanson-wright-gaussian}. Apply that bound with the matrix $\bB$ and parameter $\nu := C_1 K^2 \mu$: provided $\nu \leqslant 1/(4\opnorm{\bB})$,
\[
\expc_\bg\!\left[e^{\nu(\bg'\bB\bg - \tr(\bB))}\right] \leqslant e^{2\nu^2\fnorm{\bB}^2}, \qquad\text{i.e.}\qquad \expc_\bg\!\left[e^{\nu\,\bg'\bB\bg}\right] \leqslant \exp\!\left(\nu\,\tr(\bB) + 2\nu^2\fnorm{\bB}^2\right).
\]
Substituting $\nu = C_1 K^2\mu$ into~\eqref{eq:hw-mgf-reduced},
\[
\expc_\by\!\left[e^{\mu\,\by'\bB\by}\right] \leqslant \exp\!\left(C_1 K^2\mu\,\tr(\bB) + 2 C_1^2 K^4\mu^2\fnorm{\bB}^2\right),
\]
valid when $C_1 K^2\mu \leqslant 1/(4\opnorm{\bB})$, that is, for $0 \leqslant \mu \leqslant c_2/(K^2\opnorm{\bB})$ with $c_2 := 1/(4C_1)$. Since $\tr(\bB) \geqslant 0$ and $\mu \geqslant 0$, setting $C := \sqrt{2}\,C_1$ (so that $C \geqslant C_1$ and $C^2 = 2C_1^2$) yields~\eqref{eq:hw-mgf-quadratic}. \hfill$\square$

\textbf{Step 5: MGF of the decoupled bilinear form.} Fix a subset $I$. Taking $\expc_\by$ of the conditional bound~\eqref{eq:hw-cond-mgf} and writing $\bB = (\bA_0^{(I)})'\bA_0^{(I)}$ --- which is symmetric positive semidefinite with $\|\bA_0^{(I)}\by\|^2 = \by'\bB\by$ ---
\[
\expc\!\left[e^{\lambda\,\bx'\bA_0^{(I)}\by}\right] \leqslant \expc_\by\!\left[e^{\mu\,\by'\bB\by}\right], \qquad \mu := \frac{C_1 K^2\lambda^2}{2} \geqslant 0.
\]
The matrix $\bB$ satisfies $\tr(\bB) = \fnorm{\bA_0^{(I)}}^2$ and $\opnorm{\bB} = \opnorm{\bA_0^{(I)}}^2$, and, since the eigenvalues of $\bB$ are the squared singular values $\sigma_k^2$ of $\bA_0^{(I)}$,
\[
\fnorm{\bB}^2 = \sum_k \sigma_k^4 \leqslant \Big(\max_k \sigma_k^2\Big)\sum_k \sigma_k^2 = \opnorm{\bA_0^{(I)}}^2\,\fnorm{\bA_0^{(I)}}^2.
\]
The MGF lemma~\eqref{eq:hw-mgf-quadratic} applies when $\mu \leqslant c_2/(K^2\opnorm{\bB})$, i.e.\ when $C_1 K^2\lambda^2/2 \leqslant c_2/(K^2\opnorm{\bA_0^{(I)}}^2)$, equivalently $|\lambda| \leqslant c_3/(K^2\opnorm{\bA_0^{(I)}})$ with $c_3 := \sqrt{2c_2/C_1}$. Within this range it gives
\[
\expc\!\left[e^{\lambda\,\bx'\bA_0^{(I)}\by}\right] \leqslant \exp\!\left(\tfrac{1}{2}C C_1\, K^4\lambda^2\fnorm{\bA_0^{(I)}}^2 \;+\; \tfrac{1}{4}C^2 C_1^2\, K^8\lambda^4\opnorm{\bA_0^{(I)}}^2\fnorm{\bA_0^{(I)}}^2\right).
\]
In the second exponent term, factor out $K^4\lambda^2\fnorm{\bA_0^{(I)}}^2$: the remaining factor $K^4\lambda^2\opnorm{\bA_0^{(I)}}^2$ is at most $c_3^2$ throughout the admissible range $|\lambda| \leqslant c_3/(K^2\opnorm{\bA_0^{(I)}})$. Hence that term is bounded by $\tfrac{1}{4}C^2 C_1^2 c_3^2\, K^4\lambda^2\fnorm{\bA_0^{(I)}}^2$, and absorbing both terms into the universal constant $C_2 := \tfrac{1}{2}C C_1 + \tfrac{1}{4}C^2 C_1^2 c_3^2$,
\begin{equation}\label{eq:hw-bilinear-mgf}
\expc\!\left[e^{\lambda\,\bx'\bA_0^{(I)}\by}\right] \leqslant \exp\!\left(C_2 K^4\lambda^2\fnorm{\bA_0^{(I)}}^2\right) \quad \text{for } |\lambda| \leqslant \frac{c_3}{K^2\opnorm{\bA_0^{(I)}}}.
\end{equation}

\textbf{Step 6: Decoupling the MGF of $\bx'\bA_0\bx$.} Apply the decoupling lemma~\eqref{eq:hw-decoupling} with the convex function $F(u) = e^{\lambda u}$:
\[
\expc\!\left[e^{\lambda\,\bx'\bA_0\bx}\right] \leqslant \frac{1}{2^n}\sum_I \expc\!\left[e^{4\lambda\,\bx'\bA_0^{(I)}\by}\right].
\]
By the norm bounds of Step~2, every $I$ has $\opnorm{\bA_0^{(I)}} \leqslant \opnorm{\bA_0}$ and $\fnorm{\bA_0^{(I)}} \leqslant \fnorm{\bA_0}$. Hence if $|4\lambda| \leqslant c_3/(K^2\opnorm{\bA_0})$, then $|4\lambda| \leqslant c_3/(K^2\opnorm{\bA_0^{(I)}})$ for every $I$, and~\eqref{eq:hw-bilinear-mgf} (with $\lambda$ replaced by $4\lambda$) applies to every term of the average:
\[
\expc\!\left[e^{4\lambda\,\bx'\bA_0^{(I)}\by}\right] \leqslant \exp\!\left(16\,C_2 K^4\lambda^2\fnorm{\bA_0^{(I)}}^2\right) \leqslant \exp\!\left(16\,C_2 K^4\lambda^2\fnorm{\bA_0}^2\right).
\]
Every term obeys the same bound, hence so does their average:
\begin{equation}\label{eq:hw-offdiag-mgf}
\expc\!\left[e^{\lambda\,\bx'\bA_0\bx}\right] \leqslant \exp\!\left(C_3 K^4\lambda^2\fnorm{\bA_0}^2\right) \quad \text{for } |\lambda| \leqslant \frac{c_4}{K^2\opnorm{\bA_0}},
\end{equation}
where $C_3 := 16\,C_2$ and $c_4 := c_3/4$.

\textbf{Step 7: Chernoff bound for the off-diagonal part.} The bound~\eqref{eq:hw-offdiag-mgf} has exactly the form of~\eqref{eq:hw-gauss-mgf-bound} in the Gaussian case, with $2\fnorm{\bA}^2$ replaced by $C_3 K^4\fnorm{\bA_0}^2$ and the radius $1/(4\opnorm{\bA})$ replaced by $c_4/(K^2\opnorm{\bA_0})$. The constrained Chernoff optimization is therefore identical to the one in the proof of Theorem~\ref{thm:hanson-wright-gaussian}. For $t > 0$,
\[
\prb(\bx'\bA_0\bx \geqslant t) \leqslant \inf_{0 < \lambda \leqslant c_4/(K^2\opnorm{\bA_0})} \exp\!\left(-\lambda t + C_3 K^4\lambda^2\fnorm{\bA_0}^2\right),
\]
with unconstrained optimum $\lambda^* = t/(2C_3 K^4\fnorm{\bA_0}^2)$. If $\lambda^* \leqslant c_4/(K^2\opnorm{\bA_0})$, substituting $\lambda^*$ gives $\exp(-t^2/(4C_3 K^4\fnorm{\bA_0}^2))$. Otherwise $t > 2C_3 c_4 K^2\fnorm{\bA_0}^2/\opnorm{\bA_0}$; setting $\lambda = c_4/(K^2\opnorm{\bA_0})$ and using $\fnorm{\bA_0}^2/\opnorm{\bA_0}^2 \leqslant t/(2C_3 c_4 K^2\opnorm{\bA_0})$ gives $\exp(-c_4 t/(2K^2\opnorm{\bA_0}))$. Applying the same argument to $-\bA_0$ and a union bound,
\begin{equation}\label{eq:hw-offdiag-tail}
\prb(|\bx'\bA_0\bx| \geqslant t) \leqslant 2\exp\!\left[-c\min\!\left(\frac{t^2}{K^4\fnorm{\bA_0}^2},\; \frac{t}{K^2\opnorm{\bA_0}}\right)\right],
\end{equation}
with $c = \min(1/(4C_3),\, c_4/2)$.

\textbf{Step 8: Combine the diagonal and off-diagonal parts.} By the triangle inequality and the union bound,
\[
\prb\!\left(|\bx'\bA\bx - \expc[\bx'\bA\bx]| \geqslant t\right) \leqslant \prb\!\left(\Big|\sum_{i=1}^n a_{ii}Y_i\Big| \geqslant \tfrac{t}{2}\right) + \prb\!\left(|\bx'\bA_0\bx| \geqslant \tfrac{t}{2}\right).
\]
For the off-diagonal part, $\bA_0 = \bA - \bD$ satisfies $\fnorm{\bA_0} \leqslant \fnorm{\bA}$ (the supports of $\bA_0$ and $\bD$ are disjoint) and $\opnorm{\bA_0} \leqslant \opnorm{\bA} + \opnorm{\bD} \leqslant 2\opnorm{\bA}$. Substituting the diagonal bound~\eqref{eq:hw-diagonal} and the off-diagonal bound~\eqref{eq:hw-offdiag-tail}, each with $t$ replaced by $t/2$, and using $\fnorm{\bD} \leqslant \fnorm{\bA}$, $\opnorm{\bD} \leqslant \opnorm{\bA}$, both probabilities are bounded by $2\exp[-c\min(t^2/(K^4\fnorm{\bA}^2),\, t/(K^2\opnorm{\bA}))]$ for a universal $c > 0$ (the factors $1/2$ and the constant from $\opnorm{\bA_0} \leqslant 2\opnorm{\bA}$ are absorbed into $c$). Their sum has the same form with a possibly smaller universal constant, which is~\eqref{eq:hanson-wright}.
\end{proof}

\begin{remark}[Scope of the proof]
The argument is genuinely complete. The decoupling lemma is proved in full (Step~2) via the random-subset identity, Jensen's inequality, and an independent-copy substitution; the MGF lemma is proved in full (Step~4), the Hubbard--Stratonovich identity reducing the sub-Gaussian quadratic-form MGF to the Gaussian one; and the Gaussian quadratic-form MGF is computed explicitly in Theorem~\ref{thm:hanson-wright-gaussian} via the spectral decomposition. The Chernoff machinery and the sub-Gaussian Hoeffding inequality are derived elsewhere in the notes. No external black box is used. The original Hanson--Wright (1971) and Rudelson--Vershynin (2013) proofs take different routes through the moment-comparison principle for sub-Gaussian variables; the Hubbard--Stratonovich shortcut bypasses that machinery entirely.
\end{remark}

\begin{remark}
When $\bA$ is symmetric idempotent (as in the classical theory: $\bA = \bH$ or $\bA = \bM$), $\fnorm{\bA}^2 = \tr(\bA'\bA) = \tr(\bA) = \rank(\bA) = r$, and $\opnorm{\bA} = 1$. The Hanson--Wright inequality then gives
\[
\prb\!\left(|\bx'\bA\bx - \expc[\bx'\bA\bx]| \geqslant t\right) \leqslant 2\exp\!\left[-c\min\!\left(\frac{t^2}{K^4 r},\; \frac{t}{K^2}\right)\right].
\]
Compare this with the exact chi-square distribution $\bz'\bA\bz \sim \chi^2_r$ under Gaussianity (Theorem~4.3 in the classical notes). The Hanson--Wright bound gives the correct ``scale'' $r$ (in the Gaussian regime, $\chi^2_r$ has variance $2r$) and extends to sub-Gaussian errors with no distributional assumption.
\end{remark}

%=============================================================================
\section{Sub-Exponential Variables and Bernstein's Inequality}\label{sec:subexp}
%=============================================================================

\subsection{Definition and Basic Properties}

\begin{definition}[Sub-Exponential Random Variable]\label{def:subexponential}
A random variable $X$ is called \textbf{sub-exponential} if there exist constants $K_1, K_2 > 0$ such that
\begin{equation}\label{eq:subexp-mgf}
\expc[e^{tX}] \leqslant e^{K_1^2 t^2/2} \quad \text{for all } |t| \leqslant 1/K_2.
\end{equation}
The \textbf{sub-exponential norm} (or $\psi_1$-norm) is
\begin{equation}\label{eq:psi1}
\|X\|_{\psi_1} = \inf\!\left\{t > 0 : \expc\!\left[e^{|X|/t}\right] \leqslant 2\right\}.
\end{equation}
\end{definition}

\begin{proposition}[Squares of Sub-Gaussians Are Sub-Exponential]\label{prop:sg-square}
If $X$ is sub-Gaussian, then $X^2$ is sub-exponential with $\|X^2\|_{\psi_1} \leqslant C\|X\|_{\psi_2}^2$.
\end{proposition}

\begin{proof}
We have $\expc[e^{X^2/t}] = \expc[e^{|X|^2/t}]$. By the definition of $\|X\|_{\psi_2}$, $\expc[e^{X^2/\|X\|_{\psi_2}^2}] \leqslant 2$. Setting $t = \|X\|_{\psi_2}^2$, we get $\expc[e^{|X^2|/t}] \leqslant 2$. Hence $\|X^2\|_{\psi_1} \leqslant \|X\|_{\psi_2}^2$.
\end{proof}

\begin{example}\label{ex:subexp}
\leavevmode
\begin{enumerate}[label=\textup{(\arabic*)}]
\item $\chi^2_1$: if $Z \sim N(0,1)$, then $Z^2$ is sub-exponential by Proposition~\ref{prop:sg-square}. More generally, $\chi^2_n = \sum_{i=1}^n Z_i^2$ is sub-exponential.
\item $\mathrm{Exp}(\lambda)$: if $X \sim \mathrm{Exp}(\lambda)$ (rate $\lambda$), then $\expc[e^{X/t}] = \lambda/(\lambda - 1/t) = 2$ at $t = 2/\lambda$, so $\|X\|_{\psi_1} = 2/\lambda$.
\item Products: if $X, Y$ are sub-Gaussian, then $XY$ is sub-exponential with $\|XY\|_{\psi_1} \leqslant \|X\|_{\psi_2}\,\|Y\|_{\psi_2}$.
\end{enumerate}
\end{example}

\subsection{Bernstein's Inequality}

\begin{theorem}[Bernstein's Inequality]\label{thm:bernstein}
Let $X_1, \ldots, X_n$ be independent, mean-zero, sub-exponential random variables. Then for all $t \geqslant 0$:
\begin{equation}\label{eq:bernstein}
\boxed{\prb\!\left(\left|\sum_{i=1}^n X_i\right| \geqslant t\right) \leqslant 2\exp\!\left[-c\min\!\left(\frac{t^2}{\sum_{i=1}^n \|X_i\|_{\psi_1}^2},\; \frac{t}{\max_i \|X_i\|_{\psi_1}}\right)\right],}
\end{equation}
where $c > 0$ is a universal constant.
\end{theorem}

\begin{proof}
For $|\lambda| \leqslant 1/(C\max_i\|X_i\|_{\psi_1})$, each factor satisfies $\expc[e^{\lambda X_i}] \leqslant e^{C'^2\|X_i\|_{\psi_1}^2\lambda^2/2}$ by the sub-exponential MGF condition. The Chernoff argument gives:
\[
\prb\!\left(\sum X_i \geqslant t\right) \leqslant \exp\!\left(-\lambda t + \frac{C'^2\lambda^2}{2}\sum\|X_i\|_{\psi_1}^2\right).
\]
We optimize subject to the constraint $\lambda \leqslant 1/(C\max_i\|X_i\|_{\psi_1})$. The unconstrained optimum is $\lambda^* = t/(C'^2\sum\|X_i\|_{\psi_1}^2)$.

\textbf{Case 1:} $\lambda^*$ satisfies the constraint, i.e., $t \leqslant C'^2\sum\|X_i\|_{\psi_1}^2/(C\max_i\|X_i\|_{\psi_1})$. Substituting $\lambda^*$ gives the sub-Gaussian regime: $\exp(-t^2/(2C'^2\sum\|X_i\|_{\psi_1}^2))$.

\textbf{Case 2:} $\lambda^*$ violates the constraint. Set $\lambda = 1/(C\max_i\|X_i\|_{\psi_1})$:
\[
\prb\!\left(\sum X_i \geqslant t\right) \leqslant \exp\!\left(-\frac{t}{C\max_i\|X_i\|_{\psi_1}} + \frac{C'^2\sum\|X_i\|_{\psi_1}^2}{2C^2(\max_i\|X_i\|_{\psi_1})^2}\right).
\]
Since the constraint is violated, $t > C'^2\sum\|X_i\|_{\psi_1}^2/(C\max_i\|X_i\|_{\psi_1})$, so the constant (second) term obeys
\[
\frac{C'^2\sum\|X_i\|_{\psi_1}^2}{2C^2(\max_i\|X_i\|_{\psi_1})^2} = \frac{1}{2C\max_i\|X_i\|_{\psi_1}}\cdot\frac{C'^2\sum\|X_i\|_{\psi_1}^2}{C\max_i\|X_i\|_{\psi_1}} < \frac{t}{2C\max_i\|X_i\|_{\psi_1}}.
\]
Hence the exponent is at most $-\dfrac{t}{C\max_i\|X_i\|_{\psi_1}} + \dfrac{t}{2C\max_i\|X_i\|_{\psi_1}} = -\dfrac{t}{2C\max_i\|X_i\|_{\psi_1}}$, giving $\exp(-ct/\max_i\|X_i\|_{\psi_1})$ with $c = 1/(2C)$.

The minimum in~\eqref{eq:bernstein} captures both cases. The two-sided bound follows by the union bound.
\end{proof}

\begin{remark}[Comparison: Hoeffding vs.\ Bernstein]
For bounded variables $|X_i| \leqslant M$ (hence sub-Gaussian and sub-exponential), both inequalities apply. Hoeffding gives purely Gaussian tails $\sim e^{-ct^2}$. Bernstein also gives Gaussian tails for small $t$, but transitions to exponential tails $\sim e^{-ct}$ for $t \gtrsim \sum\|X_i\|_{\psi_1}^2/\max_i\|X_i\|_{\psi_1}$. Bernstein can be sharper because its Gaussian regime uses the variance $\sum\|X_i\|_{\psi_1}^2$, which can be much smaller than the range-based quantity $\sum(b_i - a_i)^2$ in Hoeffding's bound when the variables are concentrated near zero.
\end{remark}

%=============================================================================
\newpage
\part{Random Matrices: Covariance Estimation and Design Conditioning}\label{part:random-matrices}
%=============================================================================

%=============================================================================
\section{Covering Numbers and $\varepsilon$-Nets}\label{sec:covering}
%=============================================================================

\begin{definition}[Covering Number and $\varepsilon$-Net]\label{def:covering}
Let $(T, d)$ be a metric space and $\varepsilon > 0$. An \textbf{$\varepsilon$-net} of $T$ is a subset $\calN \subseteq T$ such that for every $\bv \in T$ there exists $\bu \in \calN$ with $d(\bv, \bu) \leqslant \varepsilon$. The \textbf{covering number} $N(T, d, \varepsilon)$ is the smallest cardinality of an $\varepsilon$-net of $T$.
\end{definition}

\begin{lemma}[Covering Number of the Unit Sphere]\label{lem:covering-sphere}
For every $\varepsilon \in (0, 1)$:
\begin{equation}\label{eq:covering-sphere}
N(S^{n-1}, \|\cdot\|, \varepsilon) \leqslant \left(\frac{2}{\varepsilon} + 1\right)^n.
\end{equation}
In particular, $N(S^{n-1}, \|\cdot\|, 1/4) \leqslant 9^n$.
\end{lemma}

\begin{proof}
Let $\calN$ be a maximal $\varepsilon$-separated subset of $S^{n-1}$: $\|\bu - \bv\| > \varepsilon$ for all distinct $\bu, \bv \in \calN$, and $\calN$ cannot be enlarged. Then $\calN$ is an $\varepsilon$-net (otherwise some point has distance $> \varepsilon$ from every element of $\calN$, contradicting maximality). The open balls $B(\bu, \varepsilon/2)$ for $\bu \in \calN$ are pairwise disjoint (if $\bw \in B(\bu_1, \varepsilon/2) \cap B(\bu_2, \varepsilon/2)$, then $\|\bu_1 - \bu_2\| \leqslant \|\bu_1 - \bw\| + \|\bw - \bu_2\| < \varepsilon$, contradicting separation). All such balls lie in $B(\bzero, 1 + \varepsilon/2)$. Comparing volumes ($\mathrm{vol}(B(\bzero, r)) = c_n r^n$):
\[
|\calN| \cdot (\varepsilon/2)^n \leqslant (1 + \varepsilon/2)^n \implies |\calN| \leqslant \left(\frac{2 + \varepsilon}{\varepsilon}\right)^n = \left(\frac{2}{\varepsilon} + 1\right)^n. \qedhere
\]
\end{proof}

\begin{lemma}[$\varepsilon$-Net Approximation for Operator Norms]\label{lem:eps-net-approx}
Let $\bA$ be a symmetric $n \times n$ matrix and $\calN$ a $(1/4)$-net of $S^{n-1}$. Then:
\begin{equation}\label{eq:eps-net-approx}
\opnorm{\bA} \leqslant 2\sup_{\bv \in \calN} |\bv'\bA\bv|.
\end{equation}
\end{lemma}

\begin{proof}
Since $\bA$ is symmetric, $\opnorm{\bA} = \sup_{\bv \in S^{n-1}} |\bv'\bA\bv|$. Let $\bv^*$ achieve this supremum and choose $\bv_0 \in \calN$ with $\|\bv^* - \bv_0\| \leqslant 1/4$. Then:
\begin{align*}
|\bv^{*\prime}\bA\bv^* - \bv_0'\bA\bv_0| &= |\bv^{*\prime}\bA(\bv^* - \bv_0) + (\bv^* - \bv_0)'\bA\bv_0|\\
&\leqslant \opnorm{\bA}\|\bv^* - \bv_0\| + \|\bv^* - \bv_0\|\opnorm{\bA} \leqslant 2\cdot\frac{1}{4}\cdot\opnorm{\bA} = \frac{1}{2}\opnorm{\bA}.
\end{align*}
Hence $\opnorm{\bA} = |\bv^{*\prime}\bA\bv^*| \leqslant |\bv_0'\bA\bv_0| + \frac{1}{2}\opnorm{\bA}$, which gives $\opnorm{\bA} \leqslant 2|\bv_0'\bA\bv_0| \leqslant 2\sup_{\bv \in \calN}|\bv'\bA\bv|$.
\end{proof}

%=============================================================================
\section{Covariance Estimation}\label{sec:covariance}
%=============================================================================

\subsection{Setup}

\begin{remark}
In the classical linear model, the design matrix $\bX$ is treated as fixed and the Gram matrix $\bX'\bX$ is assumed to have full rank. In the modern setting, the rows $\bx_1, \ldots, \bx_n \in \Real^p$ of $\bX$ are i.i.d.\ draws from a distribution with mean $\bzero$ and covariance $\bSigma = \expc[\bx_i\bx_i']$. The sample covariance is $\hbSigma = \frac{1}{n}\bX'\bX = \frac{1}{n}\sum_{i=1}^n \bx_i\bx_i'$, and the fundamental question is: \emph{how close is $\hbSigma$ to $\bSigma$?}
\end{remark}

\begin{definition}[Isotropic Random Vector]\label{def:isotropic}
A random vector $\bx \in \Real^p$ is \textbf{isotropic} if $\expc[\bx] = \bzero$ and $\expc[\bx\bx'] = \bI_p$.
\end{definition}

\begin{remark}
Any mean-zero random vector with invertible covariance $\bSigma$ can be made isotropic by the transformation $\bz = \bSigma^{-1/2}\bx$, since $\expc[\bz\bz'] = \bSigma^{-1/2}\bSigma\bSigma^{-1/2} = \bI_p$. Working with isotropic vectors simplifies notation; all results extend to general $\bSigma$ via this whitening transformation.
\end{remark}

\subsection{Operator Norm Bound for Covariance Estimation}

\begin{theorem}[Covariance Estimation for Sub-Gaussian Designs]\label{thm:cov-estimation}
Let $\bx_1, \ldots, \bx_n$ be independent, isotropic, sub-Gaussian random vectors in $\Real^p$ with $\|\bx_i\|_{\psi_2} \leqslant K$. Let $\hbSigma = \frac{1}{n}\sum_{i=1}^n \bx_i\bx_i'$. There exist universal constants $C, c > 0$ such that for every $\delta \in (0, 1)$, if $n \geqslant CK^4 p/\delta^2$, then with probability at least $1 - 2\exp(-c\delta^2 n/K^4)$:
\begin{equation}\label{eq:cov-estimation}
\boxed{\opnorm{\hbSigma - \bI_p} \leqslant \delta.}
\end{equation}
\end{theorem}

\begin{proof}
\textbf{Step 1: Variational representation.} Since $\hbSigma - \bI_p$ is symmetric:
\[
\opnorm{\hbSigma - \bI_p} = \sup_{\bv \in S^{p-1}} |\bv'(\hbSigma - \bI_p)\bv| = \sup_{\bv \in S^{p-1}}\left|\frac{1}{n}\sum_{i=1}^n \bigl[(\bv'\bx_i)^2 - 1\bigr]\right|.
\]

\textbf{Step 2: Pointwise concentration.} Fix $\bv \in S^{p-1}$. The random variable $Z_i = (\bv'\bx_i)^2 - 1$ has $\expc[Z_i] = 0$ (isotropic assumption). Since $\bv'\bx_i$ is sub-Gaussian with $\|\bv'\bx_i\|_{\psi_2} \leqslant K$, by Proposition~\ref{prop:sg-square}, $(\bv'\bx_i)^2$ is sub-exponential with $\|(\bv'\bx_i)^2\|_{\psi_1} \leqslant K^2$. Hence $Z_i$ is sub-exponential with $\|Z_i\|_{\psi_1} \leqslant CK^2$. By Bernstein's inequality (Theorem~\ref{thm:bernstein}):
\begin{equation}\label{eq:cov-pointwise}
\prb\!\left(\left|\frac{1}{n}\sum_{i=1}^n Z_i\right| \geqslant t\right) \leqslant 2\exp\!\left[-cn\min\!\left(\frac{t^2}{K^4},\; \frac{t}{K^2}\right)\right].
\end{equation}
For $t \leqslant 1$ (the regime of interest), $\min(t^2/K^4, t/K^2) = t^2/K^4$ when $t \leqslant K^2$, so the bound is $2\exp(-cnt^2/K^4)$.

\textbf{Step 3: Discretization.} Let $\calN$ be a $(1/4)$-net of $S^{p-1}$. By Lemma~\ref{lem:covering-sphere}, $|\calN| \leqslant 9^p$. By Lemma~\ref{lem:eps-net-approx}, $\opnorm{\hbSigma - \bI_p} \leqslant 2\sup_{\bv \in \calN} |\bv'(\hbSigma - \bI_p)\bv|$.

\textbf{Step 4: Union bound.} Set $t = \delta/2$ in~\eqref{eq:cov-pointwise}. For each $\bv \in \calN$:
\[
\prb\!\left(|\bv'(\hbSigma - \bI_p)\bv| \geqslant \delta/2\right) \leqslant 2\exp\!\left(-\frac{cn\delta^2}{4K^4}\right).
\]
Taking the union bound over $\calN$:
\begin{align*}
\prb\!\left(\opnorm{\hbSigma - \bI_p} \geqslant \delta\right) &\leqslant |\calN|\cdot 2\exp\!\left(-\frac{cn\delta^2}{4K^4}\right) \leqslant 2\cdot 9^p\cdot\exp\!\left(-\frac{cn\delta^2}{4K^4}\right)\\
&= 2\exp\!\left(p\ln 9 - \frac{cn\delta^2}{4K^4}\right).
\end{align*}
When $n \geqslant CK^4 p/\delta^2$ with $C = 8\ln 9/c$, the exponent is $\leqslant -cn\delta^2/(8K^4)$. Hence $\prb(\opnorm{\hbSigma - \bI_p} \geqslant \delta) \leqslant 2\exp(-c'\delta^2 n/K^4)$ with $c' = c/8$.
\end{proof}

%=============================================================================
\section{Eigenvalues of $\bX'\bX/n$: When Is the Design Well-Conditioned?}\label{sec:eigenvalues}
%=============================================================================

\subsection{From Covariance Estimation to Eigenvalue Bounds}

\begin{theorem}[Eigenvalue Bounds for Sub-Gaussian Designs]\label{thm:eigenvalue-bounds}
Under the conditions of Theorem~\ref{thm:cov-estimation}, if $n \geqslant CK^4 p$, then with probability at least $1 - 2\exp(-cp)$:
\begin{equation}\label{eq:eigenvalue-bounds}
\boxed{\frac{1}{2} \leqslant \lambda_{\min}\!\left(\frac{\bX'\bX}{n}\right) \leqslant \lambda_{\max}\!\left(\frac{\bX'\bX}{n}\right) \leqslant \frac{3}{2}.}
\end{equation}
Equivalently, $\frac{n}{2}\,\bI_p \preceq \bX'\bX \preceq \frac{3n}{2}\,\bI_p$.
\end{theorem}

\begin{proof}
Apply Theorem~\ref{thm:cov-estimation} with $\delta = 1/2$. The sample-size condition becomes $n \geqslant 4CK^4 p$. On the resulting high-probability event, $\opnorm{\bX'\bX/n - \bI_p} \leqslant 1/2$. Since $\opnorm{\bA} = \max(|\lambda_{\min}(\bA)|, |\lambda_{\max}(\bA)|)$ for symmetric $\bA$, every eigenvalue $\lambda_j$ of $\bX'\bX/n$ satisfies $|\lambda_j - 1| \leqslant 1/2$, giving $1/2 \leqslant \lambda_j \leqslant 3/2$. The probability bound is $1 - 2\exp(-c'n/(4K^4)) \geqslant 1 - 2\exp(-c'p)$ since $n \geqslant 4CK^4p$.
\end{proof}

\begin{remark}[Connection to the classical full-rank assumption]
Theorem~\ref{thm:eigenvalue-bounds} is the non-asymptotic replacement for the classical assumption (A4), $\rank(\bX) = k+1$. The classical assumption is qualitative (invertibility or not); the modern bound is quantitative: it says \emph{how well-conditioned} $\bX'\bX$ is. Specifically, $\kappa(\bX'\bX/n) = \lambda_{\max}/\lambda_{\min} \leqslant 3$ with high probability whenever $n \gtrsim K^4 p$. This is the sample-size condition under which OLS is numerically stable.
\end{remark}

\subsection{The Condition Number and Its Role}

\begin{definition}[Condition Number]\label{def:condition-number}
The \textbf{condition number} of a positive definite matrix $\bA$ is
\begin{equation}\label{eq:condition-number}
\kappa(\bA) = \frac{\lambda_{\max}(\bA)}{\lambda_{\min}(\bA)} = \opnorm{\bA}\cdot\opnorm{\bA^{-1}}.
\end{equation}
\end{definition}

\begin{proposition}[Condition Number of $\bX'\bX$ for General Covariance]\label{prop:kappa-general}
Let $\bx_1, \ldots, \bx_n$ be i.i.d.\ sub-Gaussian with covariance $\bSigma$ (not necessarily $\bI_p$). If $n \geqslant CK^4 p$, then with high probability:
\begin{equation}\label{eq:kappa-general}
\kappa\!\left(\frac{\bX'\bX}{n}\right) \leqslant 3\,\kappa(\bSigma).
\end{equation}
\end{proposition}

\begin{proof}
Apply Theorem~\ref{thm:eigenvalue-bounds} to the whitened vectors $\bz_i = \bSigma^{-1/2}\bx_i$ (which are isotropic sub-Gaussian). Then $\frac{1}{n}\sum \bz_i\bz_i' = \bSigma^{-1/2}\hbSigma\bSigma^{-1/2}$, and the eigenvalues of this matrix lie in $[1/2, 3/2]$. The eigenvalues of $\hbSigma = \bX'\bX/n$ are those of $\bSigma^{1/2}(\frac{1}{n}\sum\bz_i\bz_i')\bSigma^{1/2}$. By the min-max theorem:
\begin{align*}
\lambda_{\max}(\hbSigma) &\leqslant \frac{3}{2}\lambda_{\max}(\bSigma), &
\lambda_{\min}(\hbSigma) &\geqslant \frac{1}{2}\lambda_{\min}(\bSigma).
\end{align*}
Hence $\kappa(\hbSigma) \leqslant (3/2)/(1/2) \cdot \kappa(\bSigma) = 3\kappa(\bSigma)$.
\end{proof}

%=============================================================================
\section{The Proportional Regime $p/n \to \gamma$}\label{sec:proportional}
%=============================================================================

\subsection{Motivation: The Marchenko--Pastur Law}

\begin{remark}
When $p$ is fixed and $n \to \infty$, the law of large numbers ensures $\hbSigma \to \bSigma$ and the design is always well-conditioned for large $n$. But in modern applications (genomics, finance, natural language), $p$ can be comparable to or even larger than $n$. The \textbf{proportional regime} studies $p, n \to \infty$ with $p/n \to \gamma \in (0, \infty)$. The Marchenko--Pastur law describes the asymptotic eigenvalue distribution in this regime.
\end{remark}

\begin{theorem}[Marchenko--Pastur Law (Statement Only)]\label{thm:marchenko-pastur}
Let $\bx_1, \ldots, \bx_n \in \Real^p$ be i.i.d.\ with $\expc[\bx_i] = \bzero$ and $\expc[\bx_i\bx_i'] = \bI_p$, with i.i.d.\ entries having unit variance and finite fourth moment. Let $p/n \to \gamma \in (0, \infty)$. The empirical spectral distribution of $\hbSigma = \frac{1}{n}\sum \bx_i\bx_i'$ converges weakly almost surely to the \textbf{Marchenko--Pastur distribution} $F_\gamma$ with density
\begin{equation}\label{eq:mp-density}
f_\gamma(x) = \frac{1}{2\pi\gamma x}\sqrt{(\lambda_+ - x)(x - \lambda_-)}\,\mathbf{1}_{[\lambda_-,\,\lambda_+]}(x),
\end{equation}
where $\lambda_\pm = (1 \pm \sqrt\gamma)^2$. When $\gamma > 1$, the distribution additionally has a point mass $(1 - 1/\gamma)$ at zero.
\end{theorem}

\begin{remark}[Implications for $\bX'\bX$]
The Marchenko--Pastur law reveals three qualitatively different regimes:
\begin{enumerate}[label=\textup{(\arabic*)}]
\item \textbf{Classical regime} ($\gamma \to 0$, i.e., $p \ll n$): $\lambda_\pm \to 1$. The eigenvalues of $\hbSigma$ cluster around $1$. The design is well-conditioned; the classical theory applies.
\item \textbf{Proportional regime} ($\gamma \in (0, 1)$): $\lambda_- = (1 - \sqrt\gamma)^2 > 0$, $\lambda_+ = (1 + \sqrt\gamma)^2$. The design is invertible but the condition number $\kappa \approx \lambda_+/\lambda_- = ((1+\sqrt\gamma)/(1-\sqrt\gamma))^2 \to \infty$ as $\gamma \to 1^-$. This is the ``hard'' regime where classical methods still work but their performance degrades.
\item \textbf{Overparametrized regime} ($\gamma > 1$, i.e., $p > n$): $\lambda_- = 0$. The design matrix $\bX$ has a nontrivial null space, $\bX'\bX$ is singular, and the OLS estimator $(\bX'\bX)^{-1}\bX'\bY$ does not exist in the classical sense. Regularization (ridge, lasso) or the Moore--Penrose pseudoinverse is needed.
\end{enumerate}
\end{remark}

\subsection{Non-Asymptotic Bounds in the Proportional Regime}

\begin{theorem}[Extreme Eigenvalue Bounds, Proportional Regime]\label{thm:extreme-eigenvalues}
Let $\bx_1, \ldots, \bx_n$ be independent, isotropic, sub-Gaussian random vectors in $\Real^p$ with $\|\bx_i\|_{\psi_2} \leqslant K$. There exist constants $C, c > 0$ depending only on $K$ such that for every $t \geqslant 0$, with probability at least $1 - 2\exp(-ct^2)$:
\begin{equation}\label{eq:extreme-eigenvalues}
\sqrt{n} - C\sqrt{p} - t \leqslant \sigma_{\min}(\bX) \leqslant \sigma_{\max}(\bX) \leqslant \sqrt{n} + C\sqrt{p} + t.
\end{equation}
In terms of $\bX'\bX/n$:
\begin{equation}\label{eq:extreme-eigenvalues-gram}
\left(1 - C\sqrt{\frac{p}{n}} - \frac{t}{\sqrt{n}}\right)^2 \leqslant \lambda_{\min}\!\left(\frac{\bX'\bX}{n}\right) \leqslant \lambda_{\max}\!\left(\frac{\bX'\bX}{n}\right) \leqslant \left(1 + C\sqrt{\frac{p}{n}} + \frac{t}{\sqrt{n}}\right)^2.
\end{equation}
\end{theorem}

\begin{proof}
We prove the sharp bound for a standard Gaussian matrix---the prototype that fixes the constants---and then state the sub-Gaussian extension.

\medskip\noindent\textbf{Gaussian case.} Suppose first that the entries of $\bX$ are i.i.d.\ $N(0,1)$. Write the extreme singular values variationally:
\[
\sigma_{\max}(\bX) = \max_{\bu \in S^{p-1}}\max_{\bv \in S^{n-1}} \bv'\bX\bu, \qquad \sigma_{\min}(\bX) = \min_{\bu \in S^{p-1}}\max_{\bv \in S^{n-1}} \bv'\bX\bu.
\]
Consider the centred Gaussian process $X_{\bu,\bv} = \bv'\bX\bu$ and the comparison process $Y_{\bu,\bv} = \ip{\bg}{\bv} + \ip{\bh}{\bu}$, with $\bg \sim N(\bzero,\bI_n)$ and $\bh \sim N(\bzero,\bI_p)$ independent. Since $\expc[X_{\bu,\bv}X_{\bu',\bv'}] = \ip{\bu}{\bu'}\ip{\bv}{\bv'}$ and $\expc[Y_{\bu,\bv}Y_{\bu',\bv'}] = \ip{\bu}{\bu'} + \ip{\bv}{\bv'}$, the increments satisfy
\[
\expc(X_{\bu,\bv}-X_{\bu',\bv'})^2 = 2 - 2\ip{\bu}{\bu'}\ip{\bv}{\bv'}, \qquad \expc(Y_{\bu,\bv}-Y_{\bu',\bv'})^2 = \|\bu-\bu'\|^2 + \|\bv-\bv'\|^2,
\]
so that
\begin{equation}\label{eq:gordon-incr}
\expc(Y_{\bu,\bv}-Y_{\bu',\bv'})^2 - \expc(X_{\bu,\bv}-X_{\bu',\bv'})^2 = 2(1-\ip{\bu}{\bu'})(1-\ip{\bv}{\bv'}) \geqslant 0,
\end{equation}
with equality whenever $\bu = \bu'$.

\emph{Upper bound.} The inequality~\eqref{eq:gordon-incr} for \emph{all} pairs is the hypothesis of the Sudakov--Fernique inequality, giving $\expc\max_{\bu,\bv} X_{\bu,\bv} \leqslant \expc\max_{\bu,\bv} Y_{\bu,\bv}$. Hence
\[
\expc\,\sigma_{\max}(\bX) \leqslant \expc\Big[\max_{\bv}\ip{\bg}{\bv}\Big] + \expc\Big[\max_{\bu}\ip{\bh}{\bu}\Big] = \expc\|\bg\| + \expc\|\bh\| \leqslant \sqrt{n} + \sqrt{p},
\]
using $\expc\|\bg\| \leqslant (\expc\|\bg\|^2)^{1/2} = \sqrt{n}$ and $\expc\|\bh\| \leqslant \sqrt{p}$.

\emph{Lower bound.} The equality of increments for $\bu=\bu'$ together with~\eqref{eq:gordon-incr} is exactly the hypothesis of \textbf{Gordon's inequality} (the Gaussian min--max comparison theorem; Gordon, 1985): under these conditions $\expc\min_{\bu}\max_{\bv} X_{\bu,\bv} \geqslant \expc\min_{\bu}\max_{\bv} Y_{\bu,\bv}$. Here
\[
\min_{\bu}\max_{\bv} Y_{\bu,\bv} = \max_{\bv}\ip{\bg}{\bv} + \min_{\bu}\ip{\bh}{\bu} = \|\bg\| - \|\bh\|,
\]
so $\expc\,\sigma_{\min}(\bX) \geqslant \expc\|\bg\| - \expc\|\bh\| \geqslant \sqrt{n-1} - \sqrt{p}$, where $\expc\|\bg\| \geqslant \sqrt{n-1}$ follows from $(\expc\|\bg\|)^2 = \expc\|\bg\|^2 - \var\|\bg\| \geqslant n - 1$ (the Gaussian Poincar\'e inequality gives $\var\|\bg\| \leqslant 1$, as $\bg \mapsto \|\bg\|$ is $1$-Lipschitz).

\emph{Concentration.} The maps $\bX \mapsto \sigma_{\max}(\bX)$ and $\bX \mapsto \sigma_{\min}(\bX)$ are $1$-Lipschitz in the Frobenius norm (singular values are $1$-Lipschitz by Weyl's inequality, and $\fnorm{\cdot} \geqslant \opnorm{\cdot}$). Gaussian concentration of Lipschitz functions gives $\prb(|\,f(\bX) - \expc f(\bX)\,| \geqslant t) \leqslant 2e^{-t^2/2}$ for each. With the expectation bounds and a union bound, with probability at least $1 - 2e^{-t^2/2}$,
\[
\sqrt{n-1} - \sqrt{p} - t \leqslant \sigma_{\min}(\bX) \leqslant \sigma_{\max}(\bX) \leqslant \sqrt{n} + \sqrt{p} + t,
\]
which is~\eqref{eq:extreme-eigenvalues} with $C = 1$ (the lower-order gap $\sqrt{n}-\sqrt{n-1} \leqslant 1$ is absorbed into the additive term).

\medskip\noindent\textbf{Sub-Gaussian case.} For general isotropic sub-Gaussian rows with $\|\bx_i\|_{\psi_2} \leqslant K$, the same two-sided bound holds with $C$ depending only on $K$. This is the \emph{matrix deviation inequality} (Vershynin, \emph{HDP}, Theorem~9.1.1; and Bai--Yin, 1993, for the hard edge under finite fourth moments), whose proof replaces the Gaussian comparison by a generic-chaining bound on the sub-Gaussian process $\bu \mapsto \|\bX\bu\| - \sqrt{n}$ over $S^{p-1}$; this chaining argument is beyond our scope, and the Gaussian computation above is the prototype that fixes the sharp edges $(1 \pm \sqrt{p/n})^2$ in~\eqref{eq:extreme-eigenvalues-gram}.
\end{proof}

\begin{remark}
Setting $t = 0$ in~\eqref{eq:extreme-eigenvalues}, the singular values of $\bX$ concentrate around $[\sqrt{n} - C\sqrt{p},\, \sqrt{n} + C\sqrt{p}]$, and those of $\bX'\bX/n$ around $[(1 - C\sqrt{p/n})^2,\, (1 + C\sqrt{p/n})^2]$. This is the non-asymptotic counterpart of the Marchenko--Pastur support $[\lambda_-,\lambda_+] = [(1-\sqrt\gamma)^2, (1+\sqrt\gamma)^2]$.
\end{remark}

%=============================================================================
\newpage
\part{The Linear Model Redone}\label{part:linear-model}
%=============================================================================

%=============================================================================
\section{The Modern Setup}\label{sec:modern-setup}
%=============================================================================

\begin{definition}[The Linear Model, Modern Version]\label{def:modern-model}
The data follow
\begin{equation}\label{eq:modern-model}
\bY = \bX\bb + \be, \qquad \bX \in \Real^{n \times p},
\end{equation}
where $p$ replaces $k+1$ for the total number of regressors. The following assumptions are imposed:
\begin{enumerate}[label=\textup{(M\arabic*)}]
\item\label{M1} \textbf{Random design}: The rows $\bx_1, \ldots, \bx_n \in \Real^p$ of $\bX$ are i.i.d.\ random vectors with $\expc[\bx_i] = \bzero$ and $\expc[\bx_i\bx_i'] = \bSigma$, where $\bSigma \succ 0$.
\item\label{M2} \textbf{Sub-Gaussian design}: Each $\bx_i$ is sub-Gaussian with $\|\bx_i\|_{\psi_2} \leqslant K$.
\item\label{M3} \textbf{Sub-Gaussian noise}: The errors $\varepsilon_1, \ldots, \varepsilon_n$ are i.i.d., mean-zero, sub-Gaussian with $\|\varepsilon_i\|_{\psi_2} \leqslant \sigma$, and independent of $\bX$.
\item\label{M4} \textbf{Sample size}: $n \geqslant C_0 K^4 p$ for a sufficiently large universal constant $C_0$.
\end{enumerate}
\end{definition}

\begin{remark}[Comparison with the classical assumptions]
\leavevmode
\begin{enumerate}[label=\textup{(\arabic*)}]
\item \ref{M1} replaces the fixed-design assumption. The design is now stochastic.
\item \ref{M2} replaces (A4): it is the sub-Gaussian analogue of the full-rank condition. By Theorem~\ref{thm:eigenvalue-bounds}, \ref{M2} and~\ref{M4} together imply $\bX'\bX$ is invertible with high probability.
\item \ref{M3} replaces (A5): sub-Gaussian errors instead of normal errors. We write $\sigma_\varepsilon^2 = \var(\varepsilon_i)$ for the error variance; by Proposition~\ref{prop:sg-variance}, $\sigma_\varepsilon^2 \leqslant C\|\varepsilon_i\|_{\psi_2}^2 \leqslant C\sigma^2$, and throughout $\sigma$ is treated as a tight sub-Gaussian proxy for the noise scale, so that $\sigma_\varepsilon^2 \asymp \sigma^2$.
\item \ref{M4} is the modern analogue of $n > k + 1$ in (A4). It is a quantitative sample-size condition.
\end{enumerate}
\end{remark}

%=============================================================================
\section{OLS Prediction Error}\label{sec:prediction-error}
%=============================================================================

\begin{theorem}[Non-Asymptotic Bound on Prediction Error]\label{thm:prediction-error}
Under assumptions~\ref{M1}--\ref{M4}, with probability at least $1 - 2\exp(-cp)$:
\begin{equation}\label{eq:prediction-error}
\boxed{\frac{1}{n}\|\bX(\hbb - \bb)\|^2 \leqslant C\sigma^2\,\frac{p}{n}.}
\end{equation}
\end{theorem}

\begin{proof}
\textbf{Step 1: Express the prediction error.} The OLS estimator is $\hbb = (\bX'\bX)^{-1}\bX'\bY = \bb + (\bX'\bX)^{-1}\bX'\be$ (exactly as in the classical derivation, Equation~(2.6) in the classical notes). The prediction error is:
\begin{align}
\frac{1}{n}\|\bX(\hbb - \bb)\|^2 &= \frac{1}{n}\|\bX(\bX'\bX)^{-1}\bX'\be\|^2 = \frac{1}{n}\|\bH\be\|^2 = \frac{1}{n}\be'\bH\be, \label{eq:pred-err-qf}
\end{align}
where $\bH = \bX(\bX'\bX)^{-1}\bX'$ is the hat matrix.

\textbf{Step 2: Conditional expectation.} Conditional on $\bX$ (and hence on $\bH$), the errors $\varepsilon_1, \ldots, \varepsilon_n$ are independent sub-Gaussian with mean zero and $\|\varepsilon_i\|_{\psi_2} \leqslant \sigma$. By the Hanson--Wright inequality (Theorem~\ref{thm:hanson-wright}) applied conditionally on $\bX$:
\[
\prb\!\left(\left|\be'\bH\be - \expc[\be'\bH\be \mid \bX]\right| \geqslant t \;\middle|\; \bX\right) \leqslant 2\exp\!\left[-c\min\!\left(\frac{t^2}{\sigma^4\fnorm{\bH}^2},\; \frac{t}{\sigma^2\opnorm{\bH}}\right)\right].
\]
Since $\bH$ is idempotent with $\rank(\bH) = p$ (under the event that $\bX$ has full rank):
\begin{itemize}
\item $\opnorm{\bH} = 1$ (eigenvalues are $0$ and $1$).
\item $\fnorm{\bH}^2 = \tr(\bH) = p$.
\item $\expc[\be'\bH\be \mid \bX] = \sigma_\varepsilon^2\tr(\bH) = \sigma_\varepsilon^2 p$, where $\sigma_\varepsilon^2 = \var(\varepsilon_i) \leqslant C\sigma^2$.
\end{itemize}
Here $\sigma_\varepsilon^2 = \expc[\varepsilon_i^2]$ is the actual error variance. The Hanson--Wright bound gives, for $t = \delta\sigma^2 p$ with $\delta \in (0,1)$:
\[
\prb\!\left(|\be'\bH\be - \sigma_\varepsilon^2 p| \geqslant \delta\sigma^2 p \;\middle|\; \bX\right) \leqslant 2\exp\!\left[-c\min(\delta^2 p,\; \delta p)\right] \leqslant 2\exp(-c'\delta^2 p).
\]

\textbf{Step 3: Combine.} Setting $\delta = 1$ (say), the event $\be'\bH\be \leqslant 2\sigma^2 p$ holds with probability $\geqslant 1 - 2\exp(-cp)$ conditional on $\bX$, provided $\bX$ has full rank. By Theorem~\ref{thm:eigenvalue-bounds}, $\bX$ has full rank with probability $\geqslant 1 - 2\exp(-cp)$ under~\ref{M4}. By the union bound over both events:
\[
\frac{1}{n}\|\bX(\hbb - \bb)\|^2 = \frac{1}{n}\be'\bH\be \leqslant \frac{2\sigma^2 p}{n} = C\sigma^2\frac{p}{n}
\]
with probability $\geqslant 1 - 4\exp(-cp)$.
\end{proof}

\begin{remark}
The bound $\frac{1}{n}\|\bX(\hbb-\bb)\|^2 \lesssim \sigma^2 p/n$ says that the \textbf{in-sample prediction error} scales linearly with $p/n$. This is a non-asymptotic version of the well-known heuristic: ``each estimated parameter costs you $\sigma^2/n$ in prediction accuracy.'' In the classical theory, the exact result (under normality) is $\expc[\frac{1}{n}\|\bX(\hbb-\bb)\|^2] = \sigma^2 p/n$. The modern bound matches this rate but holds with high probability and under sub-Gaussian (not Gaussian) errors.
\end{remark}

%=============================================================================
\section{OLS Estimation Error}\label{sec:estimation-error}
%=============================================================================

\begin{theorem}[Non-Asymptotic Bound on Estimation Error]\label{thm:estimation-error}
Under assumptions~\ref{M1}--\ref{M4}, with probability at least $1 - C\exp(-cp)$:
\begin{equation}\label{eq:estimation-error}
\boxed{\|\hbb - \bb\|^2 \leqslant \frac{C\sigma^2 p}{n\,\lambda_{\min}(\bSigma)}.}
\end{equation}
More precisely, on the high-probability event from Theorem~\ref{thm:eigenvalue-bounds}:
\begin{equation}\label{eq:estimation-error-exact}
\|\hbb - \bb\|^2 = (\hbb - \bb)'\bI_p(\hbb - \bb) \leqslant \frac{2}{n\,\lambda_{\min}(\bSigma)}\|\bX(\hbb-\bb)\|^2 \leqslant \frac{C'\sigma^2 p}{n\,\lambda_{\min}(\bSigma)}.
\end{equation}
\end{theorem}

\begin{proof}
\textbf{Step 1: Relate estimation error to prediction error.} For any $\bv \in \Real^p$:
\[
\frac{1}{n}\|\bX\bv\|^2 = \bv'\!\left(\frac{\bX'\bX}{n}\right)\bv \geqslant \lambda_{\min}\!\left(\frac{\bX'\bX}{n}\right)\|\bv\|^2.
\]
Hence $\|\bv\|^2 \leqslant \frac{1}{\lambda_{\min}(\bX'\bX/n)}\cdot\frac{1}{n}\|\bX\bv\|^2$. Applying this to $\bv = \hbb - \bb$:
\begin{equation}\label{eq:est-pred-link}
\|\hbb - \bb\|^2 \leqslant \frac{1}{\lambda_{\min}(\bX'\bX/n)}\cdot\frac{1}{n}\|\bX(\hbb - \bb)\|^2.
\end{equation}

\textbf{Step 2: Bound $\lambda_{\min}(\bX'\bX/n)$ from below.} By Theorem~\ref{thm:eigenvalue-bounds} applied to the whitened design $\bZ = \bX\bSigma^{-1/2}$, on the high-probability event:
\[
\lambda_{\min}\!\left(\frac{\bX'\bX}{n}\right) \geqslant \frac{1}{2}\lambda_{\min}(\bSigma).
\]

\textbf{Step 3: Combine.} Substituting the prediction error bound (Theorem~\ref{thm:prediction-error}) and the eigenvalue bound into~\eqref{eq:est-pred-link}:
\[
\|\hbb - \bb\|^2 \leqslant \frac{2}{\lambda_{\min}(\bSigma)}\cdot C\sigma^2\frac{p}{n} = \frac{C'\sigma^2 p}{n\,\lambda_{\min}(\bSigma)}.
\]
The probability is at least $1 - C''\exp(-cp)$ by the union bound over the two high-probability events.
\end{proof}

\begin{remark}[Dependence on the condition number]
The estimation error bound~\eqref{eq:estimation-error} involves $1/\lambda_{\min}(\bSigma)$, which equals $\opnorm{\bSigma^{-1}}$. When the population covariance is well-conditioned ($\kappa(\bSigma)$ small), $\lambda_{\min}(\bSigma) \approx \lambda_{\max}(\bSigma) \approx \|\bSigma\|_{\mathrm{op}}$, and the bound is comparable to $\sigma^2 p/(n\opnorm{\bSigma})$. But when $\bSigma$ is ill-conditioned (e.g., near-multicollinearity), $\lambda_{\min}(\bSigma) \to 0$ and the estimation error explodes. This is the modern analogue of the classical observation that near-multicollinearity inflates $\var(\hat{\beta}_j) = \sigma^2[(\bX'\bX)^{-1}]_{jj}$.
\end{remark}

%=============================================================================
\section{Concentration of the Residual Sum of Squares}\label{sec:sse-concentration}
%=============================================================================

In the classical theory, $\SSE/\sigma^2 \sim \chi^2_{n-p}$ under normality (Theorem~4.5 in the classical notes). The following theorem is its non-asymptotic, sub-Gaussian counterpart.

\begin{theorem}[Concentration of $\SSE$]\label{thm:sse-concentration}
Under assumptions~\ref{M1}--\ref{M4}, with probability at least $1 - C\exp(-cp)$:
\begin{equation}\label{eq:sse-concentration}
\left|\frac{\SSE}{\sigma_\varepsilon^2} - (n-p)\right| \leqslant C'\sqrt{n-p},
\end{equation}
where $\sigma_\varepsilon^2 = \var(\varepsilon_i)$. Consequently, $\MSE = \SSE/(n-p)$ satisfies $|\MSE - \sigma_\varepsilon^2| \leqslant C'\sigma_\varepsilon^2/\sqrt{n-p}$.
\end{theorem}

\begin{proof}
\textbf{Step 1: Express $\SSE$ as a quadratic form.} We have $\SSE = \mathbf{e}'\mathbf{e} = (\bM\be)'(\bM\be) = \be'\bM\be$, where $\bM = \bI_n - \bH$ is the residual maker. On the event that $\bX$ has full column rank (which holds with probability $\geqslant 1 - 2\exp(-c_0 p)$ by Theorem~\ref{thm:eigenvalue-bounds}), $\bM$ is symmetric idempotent with $\tr(\bM) = n - p$, $\opnorm{\bM} = 1$, and $\fnorm{\bM}^2 = \tr(\bM) = n - p$.

\textbf{Step 2: Conditional expectation.} By the quadratic form expectation lemma (Lemma~3.7 of the classical notes), with $\expc[\be \mid \bX] = \bzero$ and $\var(\be \mid \bX) = \sigma_\varepsilon^2\bI_n$:
\[
\expc[\be'\bM\be \mid \bX] = \sigma_\varepsilon^2\tr(\bM) + \bzero'\bM\bzero = \sigma_\varepsilon^2(n-p).
\]

\textbf{Step 3: Concentration via Hanson--Wright.} Conditional on $\bX$, the vector $\be$ has independent, mean-zero, sub-Gaussian components with $\|\varepsilon_i\|_{\psi_2} \leqslant \sigma$. By the Hanson--Wright inequality (Theorem~\ref{thm:hanson-wright}):
\[
\prb\!\left(|\be'\bM\be - \sigma_\varepsilon^2(n-p)| \geqslant t \;\middle|\; \bX\right) \leqslant 2\exp\!\left[-c\min\!\left(\frac{t^2}{\sigma^4(n-p)},\; \frac{t}{\sigma^2}\right)\right].
\]
Here we used $\fnorm{\bM}^2 = n - p$ and $\opnorm{\bM} = 1$.

\textbf{Step 4: Choose $t$.} Apply the bound at the sub-Gaussian scale, $t = \delta\sigma^2(n-p)$ with $\delta \in (0,1]$. Then the first argument of the minimum is $t^2/(\sigma^4(n-p)) = \delta^2(n-p)$ and the second is $t/\sigma^2 = \delta(n-p) \geqslant \delta^2(n-p)$, so the minimum equals $\delta^2(n-p)$ and
\[
\prb\!\left(|\be'\bM\be - \sigma_\varepsilon^2(n-p)| \geqslant \delta\sigma^2(n-p) \;\middle|\; \bX\right) \leqslant 2\exp(-c\delta^2(n-p)).
\]
Choosing $\delta = C_0/\sqrt{n-p}$ gives $|\be'\bM\be - \sigma_\varepsilon^2(n-p)| \leqslant C_0\sigma^2\sqrt{n-p}$ with conditional probability $\geqslant 1 - 2\exp(-cC_0^2)$. Since the sub-Gaussian proxy is comparable to the noise standard deviation --- $\sigma_\varepsilon^2 \asymp \sigma^2$ (Proposition~\ref{prop:sg-variance} gives $\sigma_\varepsilon^2 \leqslant C\sigma^2$; we take $\sigma$ to be a tight proxy, so also $\sigma^2 \leqslant C\sigma_\varepsilon^2$) --- dividing by $\sigma_\varepsilon^2$ yields
\[
\left|\frac{\SSE}{\sigma_\varepsilon^2} - (n-p)\right| \leqslant C_0'\sqrt{n-p}.
\]

\textbf{Step 5: Remove conditioning.} By the union bound over the full-rank event and the concentration event:
\[
\prb\!\left(\left|\frac{\SSE}{\sigma_\varepsilon^2} - (n-p)\right| \leqslant C_0\sqrt{n-p}\right) \geqslant 1 - 2\exp(-c''C_0^2) - 2\exp(-c_0 p).
\]
Choosing $C_0$ large enough and noting $n - p \geqslant c_1 p$ under~\ref{M4}: the probability is $\geqslant 1 - C\exp(-cp)$. Dividing by $n - p$: $|\MSE - \sigma_\varepsilon^2| \leqslant C_0\sigma_\varepsilon^2/\sqrt{n-p}$.
\end{proof}

%=============================================================================
\section{Non-Asymptotic Confidence Intervals Without Normality}\label{sec:confidence}
%=============================================================================

\begin{theorem}[Sub-Gaussian Confidence Interval for $\beta_j$]\label{thm:ci-modern}
Under assumptions~\ref{M1}--\ref{M4}, for each $j \in \{1, \ldots, p\}$ and $\alpha \in (0,1)$: conditionally on $\bX$ having full rank, with probability at least $1 - \alpha$,
\begin{equation}\label{eq:ci-modern}
|\hat\beta_j - \beta_j| \leqslant C\sigma\sqrt{[(\bX'\bX)^{-1}]_{jj}\cdot\ln\frac{2}{\alpha}}.
\end{equation}
\end{theorem}

\begin{proof}
\textbf{Step 1: Isolate $\hat\beta_j - \beta_j$.} Let $\mathbf{e}_j \in \Real^p$ be the $j$-th standard basis vector. Then:
\[
\hat\beta_j - \beta_j = \mathbf{e}_j'(\hbb - \bb) = \mathbf{e}_j'(\bX'\bX)^{-1}\bX'\be = \bc'\be,
\]
where $\bc = \bX(\bX'\bX)^{-1}\mathbf{e}_j \in \Real^n$.

\textbf{Step 2: Conditional sub-Gaussian bound.} Fix $\bX$ (hence $\bc$ is fixed). Then $\bc'\be = \sum_{i=1}^n c_i\varepsilon_i$ is a sum of independent, mean-zero, sub-Gaussian random variables with $\|c_i\varepsilon_i\|_{\psi_2} \leqslant |c_i|\sigma$. By Theorem~\ref{thm:hoeffding-sg}:
\[
\prb\!\left(|\bc'\be| \geqslant t \;\middle|\; \bX\right) \leqslant 2\exp\!\left(-\frac{ct^2}{\sigma^2\|\bc\|^2}\right).
\]

\textbf{Step 3: Compute $\|\bc\|^2$.}
\[
\|\bc\|^2 = \bc'\bc = \mathbf{e}_j'(\bX'\bX)^{-1}\bX'\bX(\bX'\bX)^{-1}\mathbf{e}_j = \mathbf{e}_j'(\bX'\bX)^{-1}\mathbf{e}_j = [(\bX'\bX)^{-1}]_{jj}.
\]

\textbf{Step 4: Invert the tail bound.} Setting $2\exp(-ct^2/(\sigma^2[(\bX'\bX)^{-1}]_{jj})) = \alpha$ and solving:
\[
t = \sigma\sqrt{\frac{[(\bX'\bX)^{-1}]_{jj}}{c}\cdot\ln\frac{2}{\alpha}} = C'\sigma\sqrt{[(\bX'\bX)^{-1}]_{jj}\cdot\ln\frac{2}{\alpha}},
\]
where $C' = 1/\sqrt{c}$. Hence $|\hat\beta_j - \beta_j| \leqslant t$ with probability $\geqslant 1 - \alpha$, conditional on $\bX$.
\end{proof}

\begin{remark}[Comparison with the classical interval]
The classical $t$-based confidence interval (Theorem~6.1 of the classical notes) has half-width $t_{\alpha/2,\,n-p}\cdot s\sqrt{[(\bX'\bX)^{-1}]_{jj}}$, where $s = \sqrt{\MSE}$. For large $n - p$, the quantile $t_{\alpha/2,\,n-p} \approx z_{\alpha/2}$, and the Gaussian tail approximation gives $z_{\alpha/2} \approx \sqrt{2\ln(2/\alpha)}$. So the modern and classical widths agree up to a constant factor. The modern interval has two advantages: (1)~no normality assumption on $\be$, and (2)~finite-sample validity. Its disadvantage is that it uses the unknown $\sigma$ rather than the data-driven $s$. This can be remedied: by Theorem~\ref{thm:sse-concentration}, $s^2 = \MSE$ satisfies $|s^2/\sigma_\varepsilon^2 - 1| \leqslant C/\sqrt{n-p}$ with high probability, so for large $n - p$ one may replace $\sigma$ by $s\sqrt{2}$ (say) at the cost of doubling the constant and adding the failure probability of the $\MSE$ concentration event.
\end{remark}

%=============================================================================
\section{Comparison: Classical vs.\ Modern}\label{sec:comparison}
%=============================================================================

\begin{remark}[Side-by-side comparison]
The following table compares the classical and modern treatments of the linear model. The notation $p = k+1$.
\end{remark}

\medskip
\begin{center}
\renewcommand{\arraystretch}{1.4}
\begin{tabular}{@{}p{3.2cm}p{5.5cm}p{5.5cm}@{}}
\toprule
\textbf{Aspect} & \textbf{Classical Theory} & \textbf{Modern Theory} \\
\midrule
Design & Fixed ($\bX$ non-random) & Random, i.i.d.\ rows \\
Error distribution & $\be \sim N(\bzero, \sigma^2\bI_n)$ & Sub-Gaussian, $\|\varepsilon_i\|_{\psi_2} \leqslant \sigma$ \\
Invertibility & $\rank(\bX) = p$ (assumed) & $n \geqslant CK^4 p$ (implies w.h.p.) \\
Prediction error & $\expc[\|\bH\be\|^2/n] = \sigma^2 p/n$ (exact) & $\|\bH\be\|^2/n \leqslant C\sigma^2 p/n$ (w.h.p.) \\
Estimation error & $\var(\hat{\bb}) = \sigma^2(\bX'\bX)^{-1}$ & $\|\hbb - \bb\|^2 \leqslant C\sigma^2 p/(n\lambda_{\min}(\bSigma))$ (w.h.p.) \\
Distribution of $\SSE$ & $\SSE/\sigma^2 \sim \chi^2_{n-p}$ & $|\SSE/\sigma^2 - (n-p)| \leqslant C\sqrt{n-p}$ (w.h.p.) \\
Inference & $t$, $F$, $\chi^2$ pivots & Sub-Gaussian concentration (Thm.~\ref{thm:ci-modern}) \\
Regime & $p$ fixed, $n \to \infty$ & $p/n \to \gamma \in (0,1)$ allowed \\
\bottomrule
\end{tabular}
\end{center}

%=============================================================================
\section{What Happens When $p$ Is Comparable to $n$}\label{sec:p-comparable-n}
%=============================================================================

\begin{remark}
As $p/n \to \gamma \in (0,1)$, the OLS estimator remains well-defined but its behavior changes qualitatively.
\end{remark}

\begin{theorem}[Prediction and Estimation Error in the Proportional Regime]\label{thm:proportional-prediction}
Under assumptions~\ref{M1}--\ref{M3} with $\bSigma = \bI_p$ (isotropic design), conditional on $\bX$ having full column rank:
\begin{enumerate}[label=\textup{(\roman*)}]
\item\label{pp:in} \textbf{In-sample prediction error}: $\expc[\frac{1}{n}\|\bX(\hbb-\bb)\|^2 \mid \bX] = \sigma_\varepsilon^2\,p/n$, where $\sigma_\varepsilon^2 = \var(\varepsilon_i)$.
\item\label{pp:mse} \textbf{Unbiasedness of MSE}: $\expc[\MSE \mid \bX] = \sigma_\varepsilon^2$.
\item\label{pp:oos} \textbf{Out-of-sample prediction risk}: For a new $\bx_0$ drawn independently from the same distribution,
\begin{equation}\label{eq:oos-risk}
\expc[(\bx_0'(\hbb-\bb))^2 \mid \bX] = \sigma_\varepsilon^2\tr((\bX'\bX)^{-1}).
\end{equation}
As $p/n \to \gamma \in (0,1)$ with i.i.d.\ entries having finite fourth moment, $\frac{1}{p}\tr((\bX'\bX/n)^{-1}) \to \frac{1}{1-\gamma}$ almost surely by the Marchenko--Pastur law (Theorem~\ref{thm:marchenko-pastur}), giving $\expc[(\bx_0'(\hbb-\bb))^2 \mid \bX] \to \sigma_\varepsilon^2\cdot\frac{\gamma}{1-\gamma}$.
\end{enumerate}
\end{theorem}

\begin{proof}
\textbf{(i).} Conditional on $\bX$: $\|\bX(\hbb-\bb)\|^2 = \be'\bH\be$ (from~\eqref{eq:pred-err-qf}). By the quadratic form expectation lemma (Lemma~3.7 of the classical notes) with $\expc[\be \mid \bX] = \bzero$ and $\var(\be \mid \bX) = \sigma_\varepsilon^2\bI_n$:
\[
\expc[\be'\bH\be \mid \bX] = \sigma_\varepsilon^2\tr(\bH) + \bzero'\bH\bzero = \sigma_\varepsilon^2 p.
\]
Dividing by $n$ gives~\ref{pp:in}.

\textbf{(ii).} Similarly, $\SSE = \be'\bM\be$ where $\bM = \bI - \bH$ has $\tr(\bM) = n - p$. So $\expc[\SSE \mid \bX] = \sigma_\varepsilon^2(n-p)$ and $\expc[\MSE \mid \bX] = \sigma_\varepsilon^2$.

\textbf{(iii).} Since $\bx_0$ is independent of $(\bX, \be)$ with $\expc[\bx_0\bx_0'] = \bI_p$:
\[
\expc[(\bx_0'(\hbb-\bb))^2 \mid \bX, \be] = (\hbb-\bb)'\expc[\bx_0\bx_0'](\hbb-\bb) = \|\hbb-\bb\|^2.
\]
Using $\hbb - \bb = (\bX'\bX)^{-1}\bX'\be$:
\[
\|\hbb-\bb\|^2 = \be'\bX(\bX'\bX)^{-2}\bX'\be.
\]
Let $\bC = \bX(\bX'\bX)^{-2}\bX'$. By the quadratic form lemma, $\expc[\be'\bC\be \mid \bX] = \sigma_\varepsilon^2\tr(\bC)$, and using the cyclic property of the trace:
\[
\tr(\bC) = \tr(\bX(\bX'\bX)^{-2}\bX') = \tr((\bX'\bX)^{-2}\bX'\bX) = \tr((\bX'\bX)^{-1}).
\]
Hence $\expc[(\bx_0'(\hbb-\bb))^2 \mid \bX] = \sigma_\varepsilon^2\tr((\bX'\bX)^{-1})$.

\textbf{Asymptotic statement.} The eigenvalues of $(\bX'\bX)^{-1}$ are the reciprocals of those of $\bX'\bX$, so
\[
\tr((\bX'\bX)^{-1}) = \sum_{j=1}^p \frac{1}{\lambda_j(\bX'\bX)} = \frac{1}{n}\sum_{j=1}^p \frac{1}{\lambda_j(\bX'\bX/n)}.
\]
We wish to conclude
\begin{equation}\label{eq:mp-trace-limit}
\frac{1}{p}\sum_{j=1}^p\frac{1}{\lambda_j(\bX'\bX/n)} \xrightarrow{\text{a.s.}} \int\frac{1}{x}\,\dd F_\gamma(x) = \frac{1}{1-\gamma},
\end{equation}
where the moment $\int x^{-1}\,\dd F_\gamma(x) = 1/(1-\gamma)$ for $\gamma \in (0,1)$ is a standard computation via the Stieltjes transform.

\medskip\noindent\emph{Justification of~\eqref{eq:mp-trace-limit}.} The Marchenko--Pastur law gives weak convergence of the empirical spectral distribution $F_p$ of $\bX'\bX/n$ to $F_\gamma$. Since the test function $g(x) = 1/x$ has a pole at zero, this is \emph{not} immediate from weak convergence and requires uniform truncation of the support away from zero. Under the sub-Gaussian assumption~\ref{M2}, Theorem~\ref{thm:extreme-eigenvalues} with $t = O(\sqrt{n})$ gives, for every $\gamma \in (0,1)$ and all sufficiently large $n,p$ with $p/n \to \gamma$,
\[
\lambda_{\min}(\bX'\bX/n) \geqslant (1 - C\sqrt{\gamma})^2 - o(1) =: \underline{\lambda} > 0
\]
on an event of probability $1 - 2\exp(-c n)$, which is summable in $n$; by Borel--Cantelli, $\lambda_{\min}(\bX'\bX/n) \geqslant \underline{\lambda}/2$ for all $n$ large enough almost surely.\footnote{Under the weaker moment assumption of MP (i.i.d.\ entries with finite fourth moment), the corresponding hard-edge result is the Bai--Yin law (1993). Under~\ref{M2}, the non-asymptotic bound of Theorem~\ref{thm:extreme-eigenvalues} suffices.} On this event, all eigenvalues $\lambda_j(\bX'\bX/n)$ lie in $[\underline{\lambda}/2, M]$ for a deterministic upper bound $M$ (similarly from Theorem~\ref{thm:extreme-eigenvalues}), and the truncated test function $g(x) = 1/x$ is bounded and continuous on this interval. Weak convergence then yields~\eqref{eq:mp-trace-limit}.

\medskip\noindent Combining,
\[
\sigma_\varepsilon^2\tr((\bX'\bX)^{-1}) = \sigma_\varepsilon^2 \cdot \frac{p}{n}\cdot\frac{1}{p}\sum_{j=1}^p\frac{1}{\lambda_j(\bX'\bX/n)} \xrightarrow{\text{a.s.}} \sigma_\varepsilon^2\cdot\gamma\cdot\frac{1}{1-\gamma} = \sigma_\varepsilon^2\cdot\frac{\gamma}{1-\gamma}. \qedhere
\]
\end{proof}

\begin{corollary}[Estimation Error Diverges as $\gamma \to 1^-$]\label{cor:estimation-diverges}
Under the conditions of Theorem~\ref{thm:proportional-prediction}, $\expc[\|\hbb - \bb\|^2 \mid \bX] = \sigma_\varepsilon^2\tr((\bX'\bX)^{-1}) \to \infty$ as $p/n \to 1^-$, since $\frac{1}{1-\gamma} \to \infty$.
\end{corollary}

\begin{remark}[Diagnosis and cure]
Corollary~\ref{cor:estimation-diverges} is the diagnosis: as the design approaches singularity ($\gamma \to 1^-$), the smallest eigenvalue of $\bX'\bX/n$ approaches zero, so $\opnorm{(\bX'\bX)^{-1}} \to \infty$ and OLS estimation error blows up. The standard cure is \textbf{ridge regression}: replace $\bX'\bX$ by $\bX'\bX + \lambda\bI_p$ for some $\lambda > 0$. This shifts every eigenvalue by $\lambda$, hence $\opnorm{(\bX'\bX + \lambda\bI_p)^{-1}} \leqslant 1/\lambda$ regardless of how singular $\bX'\bX$ is. The cost is a bias of order $\lambda\opnorm{\bb}$; the optimal $\lambda$ balances bias against variance and depends on $\gamma$, $\opnorm{\bb}$, and $\sigma_\varepsilon^2$. The bias--variance decomposition is the subject of Exercise~\ref{ex:mC2}.
\end{remark}

\begin{remark}[The overparametrized regime $p > n$]
When $p > n$, $\bX'\bX$ is singular and the classical OLS estimator is undefined. The minimum-norm interpolator $\hbb_{\mathrm{mn}} = \bX'(\bX\bX')^{-1}\bY$ can still be computed, and---surprisingly---its prediction risk can decrease as $p/n$ increases beyond $1$. This ``double descent'' phenomenon is the starting point of modern interpolation theory (Hastie, Montanari, Rosset, and Tibshirani, 2022), but lies beyond the scope of these notes.
\end{remark}

%=============================================================================
\section{Exercises}\label{sec:modern-exercises}
%=============================================================================

\setcounter{exercise}{0}

\subsection*{Tier A: Reproduce the Derivation}

\begin{exercise}\label{ex:mA1}
Let $X$ be a random variable with $\expc[X] = 0$ and $a \leqslant X \leqslant b$. Prove Hoeffding's lemma: $\expc[e^{tX}] \leqslant e^{t^2(b-a)^2/8}$.
\end{exercise}

\begin{exercise}\label{ex:mA2}
Let $X_1, \ldots, X_n$ be independent with $\expc[X_i] = 0$ and $a_i \leqslant X_i \leqslant b_i$. Prove Hoeffding's inequality: $\prb(|\sum X_i| \geqslant t) \leqslant 2\exp(-2t^2/\sum(b_i - a_i)^2)$.
\end{exercise}

\begin{exercise}\label{ex:mA3}
Prove that if $X$ is sub-Gaussian, then $X^2$ is sub-exponential with $\|X^2\|_{\psi_1} \leqslant C\|X\|_{\psi_2}^2$.
\end{exercise}

\begin{exercise}\label{ex:mA4}
Let $\bx_1, \ldots, \bx_n$ be i.i.d.\ isotropic sub-Gaussian in $\Real^p$. Carry out the pointwise Bernstein argument (Step~2 of Theorem~\ref{thm:cov-estimation}) to show that for a fixed $\bv \in S^{p-1}$:
\[
\prb\!\left(\left|\frac{1}{n}\sum_{i=1}^n [(\bv'\bx_i)^2 - 1]\right| \geqslant t\right) \leqslant 2\exp\!\left[-cn\min\!\left(\frac{t^2}{K^4}, \frac{t}{K^2}\right)\right].
\]
\end{exercise}

\begin{exercise}\label{ex:mA5}
Derive the prediction error identity $\frac{1}{n}\|\bX(\hbb - \bb)\|^2 = \frac{1}{n}\be'\bH\be$ and compute $\expc[\be'\bH\be \mid \bX]$ under sub-Gaussian errors.
\end{exercise}

\subsection*{Tier B: Apply the Theory}

\begin{exercise}\label{ex:mB1}
Let $X_1, \ldots, X_n$ be i.i.d.\ $\mathrm{Bernoulli}(p)$ with $p \in (0,1)$. Use Hoeffding's inequality to find the smallest $n$ guaranteeing $\prb(|\bar{X} - p| \geqslant 0.01) \leqslant 0.05$.
\end{exercise}

\begin{exercise}\label{ex:mB2}
Suppose $\bx_i \sim N(\bzero, \bSigma)$ with $\bSigma = \diag(1, 2, 3, \ldots, p)$. (a)~Compute $\kappa(\bSigma)$. (b)~How large must $n$ be (in terms of $p$) to guarantee $\opnorm{\hbSigma - \bSigma}/\opnorm{\bSigma} \leqslant 1/2$ with probability $\geqslant 0.99$? (c)~Compare the estimation error bound from Theorem~\ref{thm:estimation-error} with the classical variance formula $\sigma^2[(\bX'\bX)^{-1}]_{11}$.
\end{exercise}

\begin{exercise}\label{ex:mB3}
In the Marchenko--Pastur law with $\gamma = 0.5$, compute $\lambda_- = (1-\sqrt{0.5})^2$ and $\lambda_+ = (1+\sqrt{0.5})^2$. What is the asymptotic condition number of $\bX'\bX/n$? If the true error variance is $\sigma^2 = 1$, estimate $\|\hbb - \bb\|^2$ using Corollary~\ref{cor:estimation-diverges}.
\end{exercise}

\begin{exercise}\label{ex:mB4}
A design matrix $\bX \in \Real^{100 \times 10}$ has i.i.d.\ Rademacher entries ($\pm 1$ with equal probability). (a)~Verify that each row is sub-Gaussian and compute $K$. (b)~Check the sample-size condition $n \geqslant CK^4 p$. (c)~Give a high-probability interval for $\lambda_{\min}(\bX'\bX/n)$ and $\lambda_{\max}(\bX'\bX/n)$ using Theorem~\ref{thm:extreme-eigenvalues}.
\end{exercise}

\subsection*{Tier C: Extend and Connect}

\begin{exercise}\label{ex:mC1}
\textbf{(Concentration of the MSE.)} Under~\ref{M1}--\ref{M4}, the classical notes show $\expc[\MSE] = \sigma^2$ under normality. Use the Hanson--Wright inequality to prove that $|\SSE - \sigma^2(n-p)| \leqslant C\sigma^2\sqrt{n-p}$ with high probability under sub-Gaussian errors. Conclude that $\MSE = \SSE/(n-p)$ concentrates around $\sigma^2$.
\end{exercise}

\begin{exercise}\label{ex:mC2}
\textbf{(Ridge regression in the proportional regime.)} Consider the ridge estimator $\hbb_\lambda = (\bX'\bX + \lambda\bI_p)^{-1}\bX'\bY$. (a)~Show that $\hbb_\lambda - \bb = (\bX'\bX + \lambda\bI_p)^{-1}(\bX'\be - \lambda\bb)$. (b)~Decompose $\|\hbb_\lambda - \bb\|^2$ into a bias term (depending on $\lambda\bb$) and a variance term (depending on $\be$). (c)~Why does ridge help when $p/n \to 1^-$?
\end{exercise}

\begin{exercise}\label{ex:mC3}
\textbf{(Sub-Gaussian tail optimality.)} Show that the Hoeffding bound for $\bar{X}$ (Corollary~\ref{cor:hoeffding-iid}) cannot be improved beyond constant factors for Rademacher variables: explicitly compute $\prb(\sum_{i=1}^n \varepsilon_i \geqslant t)$ for $\varepsilon_i$ i.i.d.\ uniform on $\{-1, +1\}$ via the MGF, and compare with~\eqref{eq:hoeffding}.
\end{exercise}

\begin{exercise}\label{ex:mC4}
\textbf{(From concentration to confidence intervals.)} Under~\ref{M1}--\ref{M4}, construct a non-asymptotic confidence interval for a single coefficient $\beta_j$ \emph{without assuming normality}. Specifically: (a)~Show $\hat\beta_j - \beta_j = \mathbf{e}_j'(\bX'\bX)^{-1}\bX'\be$, where $\mathbf{e}_j$ is the $j$-th standard basis vector. (b)~Conditionally on $\bX$, bound $|\hat\beta_j - \beta_j|$ using the sub-Gaussian tail of $\bc'\be$ for $\bc = (\bX'\bX)^{-1}\bX'\mathbf{e}_j$. (c)~Show this yields a confidence interval of width $O(\sigma\sqrt{[(\bX'\bX)^{-1}]_{jj}\log(1/\alpha)})$ at level $1-\alpha$. Compare with the classical interval width $t_{\alpha/2, n-p}\cdot s\sqrt{[(\bX'\bX)^{-1}]_{jj}}$.
\end{exercise}

%=============================================================================
\section{Solutions to Exercises}\label{sec:modern-solutions}
%=============================================================================

\subsection*{Solutions to Tier A}

\textbf{Solution to Exercise~\ref{ex:mA1}.} See the proof of Lemma~\ref{lem:hoeffding-lemma}.

\textbf{Solution to Exercise~\ref{ex:mA2}.} See the proof of Theorem~\ref{thm:hoeffding}.

\textbf{Solution to Exercise~\ref{ex:mA3}.} See the proof of Proposition~\ref{prop:sg-square}.

\textbf{Solution to Exercise~\ref{ex:mA4}.} The random variable $Z_i = (\bv'\bx_i)^2 - 1$ satisfies $\expc[Z_i] = 0$ (isotropic assumption) and is sub-exponential with $\|Z_i\|_{\psi_1} \leqslant CK^2$ (since $\bv'\bx_i$ is sub-Gaussian with norm $\leqslant K$ and squares of sub-Gaussians are sub-exponential by Proposition~\ref{prop:sg-square}). By Bernstein's inequality (Theorem~\ref{thm:bernstein}) applied to $\sum_{i=1}^n Z_i$:
\[
\prb\!\left(\left|\sum_{i=1}^n Z_i\right| \geqslant nt\right) \leqslant 2\exp\!\left[-c\min\!\left(\frac{n^2t^2}{nC^2K^4},\; \frac{nt}{CK^2}\right)\right] = 2\exp\!\left[-cn\min\!\left(\frac{t^2}{K^4},\; \frac{t}{K^2}\right)\right].
\]

\textbf{Solution to Exercise~\ref{ex:mA5}.} Substituting $\hbb - \bb = (\bX'\bX)^{-1}\bX'\be$:
\begin{align*}
\bX(\hbb - \bb) &= \bX(\bX'\bX)^{-1}\bX'\be = \bH\be.
\end{align*}
Hence $\frac{1}{n}\|\bX(\hbb-\bb)\|^2 = \frac{1}{n}\|\bH\be\|^2 = \frac{1}{n}(\bH\be)'\bH\be = \frac{1}{n}\be'\bH'\bH\be = \frac{1}{n}\be'\bH\be$ (using $\bH' = \bH$ and $\bH^2 = \bH$). For the conditional expectation: since $\varepsilon_1, \ldots, \varepsilon_n$ are independent with $\expc[\varepsilon_i] = 0$ and $\expc[\varepsilon_i^2] = \sigma_\varepsilon^2$, we have $\expc[\be\be' \mid \bX] = \sigma_\varepsilon^2\bI_n$. By Lemma~3.7 of the classical notes (expected value of quadratic form, with $\boldsymbol\mu = \bzero$):
\[
\expc[\be'\bH\be \mid \bX] = \tr(\bH\cdot\sigma_\varepsilon^2\bI_n) = \sigma_\varepsilon^2\tr(\bH) = \sigma_\varepsilon^2 p.
\]
Dividing by $n$: $\expc[\frac{1}{n}\|\bX(\hbb-\bb)\|^2 \mid \bX] = \sigma_\varepsilon^2 p/n$.

\subsection*{Solutions to Tier B}

\textbf{Solution to Exercise~\ref{ex:mB1}.} Each $X_i \in [0, 1]$, so $b_i - a_i = 1$. By Corollary~\ref{cor:hoeffding-iid}: $\prb(|\bar{X} - p| \geqslant 0.01) \leqslant 2\exp(-2n(0.01)^2) = 2e^{-n/5000}$. Setting $2e^{-n/5000} \leqslant 0.05$ gives $n \geqslant 5000\ln(40) \approx 18\,444$.

\textbf{Solution to Exercise~\ref{ex:mB2}.} (a) $\kappa(\bSigma) = \lambda_{\max}/\lambda_{\min} = p/1 = p$. (b) The whitened vectors $\bz_i = \bSigma^{-1/2}\bx_i$ are isotropic Gaussian with $\|\bz_i\|_{\psi_2} \leqslant C$. The condition $\opnorm{\hbSigma - \bSigma}/\opnorm{\bSigma} \leqslant 1/2$ reduces (after whitening) to $\opnorm{\hat{\mathbf{I}} - \bI_p} \leqslant 1/2$ for the whitened sample covariance. By Theorem~\ref{thm:cov-estimation} with $\delta = 1/2$, we need $n \geqslant 4Cp$. For probability $0.99$, we additionally need $2\exp(-cp) \leqslant 0.01$, i.e., $p \geqslant \frac{1}{c}\ln 200$, which is mild.

(c) The classical formula gives $\var(\hat\beta_1) = \sigma^2[(\bX'\bX)^{-1}]_{11}$, which depends on the specific realization of $\bX$. For the diagonal covariance, $\expc[(\bX'\bX)^{-1}] \approx \frac{1}{n}\bSigma^{-1}$ for large $n$, so $\var(\hat\beta_1) \approx \sigma^2/(n\cdot 1) = \sigma^2/n$. The modern bound gives $\|\hbb - \bb\|^2 \leqslant C\sigma^2 p/(n\lambda_{\min}(\bSigma)) = C\sigma^2 p/n$, which is larger by a factor $p$ (but bounds the total $\ell^2$ error, not just one coordinate).

\textbf{Solution to Exercise~\ref{ex:mB3}.} $\lambda_- = (1 - 1/\sqrt{2})^2 = (2 - \sqrt{2})^2/4 \approx 0.0858$. $\lambda_+ = (1 + 1/\sqrt{2})^2 = (2 + \sqrt{2})^2/4 \approx 2.914$. Asymptotic condition number: $\kappa \approx \lambda_+/\lambda_- \approx 33.97$. For the estimation error in the isotropic case ($\bSigma = \bI_p$, $\sigma^2 = 1$), Theorem~\ref{thm:proportional-prediction}(iii) gives, as $p/n \to \gamma = 0.5$,
\[
\expc[\|\hbb-\bb\|^2 \mid \bX] \to \sigma_\varepsilon^2\cdot\frac{\gamma}{1-\gamma} = 1\cdot\frac{0.5}{0.5} = 1.
\]

\textbf{Solution to Exercise~\ref{ex:mB4}.} (a) Each entry $X_{ij}$ is Rademacher: $\|X_{ij}\|_{\psi_2} = C$ (by Example~\ref{ex:subgaussian}(2)). A row $\bx_i = (X_{i1}, \ldots, X_{ip})'$ has $\|\ip{\bx_i}{\bv}\|_{\psi_2} = \|\sum_j v_j X_{ij}\|_{\psi_2} \leqslant C\|\bv\|$ by Proposition~\ref{prop:sg-closure}. So $K = C$. (b) $n = 100$, $p = 10$, $K^4 p = C^4 \cdot 10$. For reasonable $C$, $n = 100 \gg C^4 \cdot 10$, so the condition is satisfied. (c) By~\eqref{eq:extreme-eigenvalues-gram} with $t = 0$: $\lambda_{\min}(\bX'\bX/n) \gtrsim (1 - C'\sqrt{10/100})^2 = (1 - C'\sqrt{0.1})^2$ and $\lambda_{\max}(\bX'\bX/n) \lesssim (1 + C'\sqrt{0.1})^2$. With $C' \approx 1$: $\lambda_{\min} \gtrsim (1 - 0.316)^2 \approx 0.467$, $\lambda_{\max} \lesssim (1 + 0.316)^2 \approx 1.732$.

\subsection*{Solutions to Tier C}

\textbf{Solution to Exercise~\ref{ex:mC1}.} We have $\SSE = \be'\bM\be$ where $\bM = \bI_n - \bH$ is symmetric idempotent with $\tr(\bM) = n - p$ and $\opnorm{\bM} = 1$, $\fnorm{\bM}^2 = n - p$. By Hanson--Wright (Theorem~\ref{thm:hanson-wright}), conditionally on $\bX$:
\[
\prb\!\left(|\SSE - \sigma_\varepsilon^2(n-p)| \geqslant t \;\middle|\; \bX\right) \leqslant 2\exp\!\left[-c\min\!\left(\frac{t^2}{\sigma^4(n-p)},\; \frac{t}{\sigma^2}\right)\right].
\]
Setting $t = \delta\sigma^2(n-p)$: $\prb(|\SSE - \sigma_\varepsilon^2(n-p)| \geqslant \delta\sigma^2(n-p)) \leqslant 2\exp(-c\min(\delta^2(n-p), \delta(n-p))) \leqslant 2\exp(-c'\delta^2(n-p))$. With $\delta = C/\sqrt{n-p}$: $|\SSE - \sigma_\varepsilon^2(n-p)| \leqslant C\sigma^2\sqrt{n-p}$ with probability $\geqslant 1 - 2e^{-c'C^2}$. Hence $|\MSE - \sigma_\varepsilon^2| = |\SSE/(n-p) - \sigma_\varepsilon^2| \leqslant C\sigma^2/\sqrt{n-p} \to 0$.

\textbf{Solution to Exercise~\ref{ex:mC2}.}

\emph{Part (a): the residual identity.} Starting from $\hbb_\lambda = (\bX'\bX + \lambda\bI)^{-1}\bX'\bY = (\bX'\bX + \lambda\bI)^{-1}\bX'(\bX\bb + \be)$, write $\bX'\bX = (\bX'\bX + \lambda\bI) - \lambda\bI$. Then
\[
(\bX'\bX + \lambda\bI)^{-1}\bX'\bX\bb = \bb - \lambda(\bX'\bX + \lambda\bI)^{-1}\bb,
\]
so $\hbb_\lambda - \bb = (\bX'\bX + \lambda\bI)^{-1}\bX'\be - \lambda(\bX'\bX + \lambda\bI)^{-1}\bb$.

\emph{Part (b): bias--variance decomposition.} Set $\bR_\lambda := (\bX'\bX + \lambda\bI)^{-1}$. Then $\hbb_\lambda - \bb = \bR_\lambda\bX'\be - \lambda\bR_\lambda\bb$, and
\[
\|\hbb_\lambda - \bb\|^2 = \lambda^2\bb'\bR_\lambda^2\bb \;-\; 2\lambda\,\bb'\bR_\lambda^2\bX'\be \;+\; \be'\bX\bR_\lambda^2\bX'\be.
\]
Taking conditional expectation given $\bX$: the cross term vanishes since $\expc[\be \mid \bX] = \bzero$, so
\begin{equation}\label{eq:ridge-decomp}
\expc[\|\hbb_\lambda - \bb\|^2 \mid \bX] = \underbrace{\lambda^2\bb'\bR_\lambda^2\bb}_{\mathrm{bias}^2(\lambda)} + \underbrace{\sigma_\varepsilon^2\,\tr(\bX\bR_\lambda^2\bX')}_{\mathrm{variance}(\lambda)},
\end{equation}
where the second equality uses $\expc[\be\be' \mid \bX] = \sigma_\varepsilon^2\bI_n$ and the cyclic property of the trace.

\emph{Spectral form.} Let $\bX'\bX = \bU\diag(\mu_1, \ldots, \mu_p)\bU'$ be the spectral decomposition, and write $\widetilde\bb := \bU'\bb$ in the eigenbasis. Then $\bR_\lambda = \bU\diag(1/(\mu_j+\lambda))\bU'$, and:
\begin{align}
\mathrm{bias}^2(\lambda) &= \lambda^2 \sum_{j=1}^p \frac{\widetilde b_j^2}{(\mu_j + \lambda)^2}, \label{eq:ridge-bias}\\
\mathrm{variance}(\lambda) &= \sigma_\varepsilon^2 \sum_{j=1}^p \frac{\mu_j}{(\mu_j + \lambda)^2}, \label{eq:ridge-variance}
\end{align}
the last using $\tr(\bX\bR_\lambda^2\bX') = \tr(\bR_\lambda^2 \bX'\bX) = \sum_j \mu_j/(\mu_j+\lambda)^2$.

\emph{Two uniform bounds.} For each $j$, the function $\lambda \mapsto \lambda^2/(\mu_j+\lambda)^2$ is monotone increasing in $\lambda$ and bounded by $1$, so
\begin{equation}\label{eq:ridge-bias-bound}
\mathrm{bias}^2(\lambda) \leqslant \opnorm{\bb}^2 \quad \text{for every } \lambda \geqslant 0.
\end{equation}
For each $j$, the function $\mu \mapsto \mu/(\mu + \lambda)^2$ is maximized at $\mu = \lambda$ with maximum value $1/(4\lambda)$, so $\mu_j/(\mu_j + \lambda)^2 \leqslant 1/(4\lambda)$, giving
\begin{equation}\label{eq:ridge-variance-bound}
\mathrm{variance}(\lambda) \leqslant \frac{\sigma_\varepsilon^2 p}{4\lambda} \quad \text{for every } \lambda > 0.
\end{equation}

\emph{The crucial contrast with OLS.} Bounds~\eqref{eq:ridge-bias-bound}--\eqref{eq:ridge-variance-bound} \emph{do not depend on $\lambda_{\min}(\bX'\bX)$.} Combining,
\[
\expc[\|\hbb_\lambda - \bb\|^2 \mid \bX] \leqslant \opnorm{\bb}^2 + \frac{\sigma_\varepsilon^2 p}{4\lambda}
\]
for every $\lambda > 0$, regardless of how singular $\bX'\bX$ is. Setting $\lambda \to \infty$ kills the variance; setting $\lambda \to 0^+$ recovers the unbounded OLS variance (cf.\ Corollary~\ref{cor:estimation-diverges}). In contrast, OLS has $\expc[\|\hbb - \bb\|^2 \mid \bX] = \sigma_\varepsilon^2 \tr((\bX'\bX)^{-1}) = \sigma_\varepsilon^2 \sum_j 1/\mu_j$, which diverges as $\mu_p = \lambda_{\min}(\bX'\bX) \to 0$.

\emph{Part (c): why ridge helps as $\gamma \to 1^-$.} In OLS, the variance is dominated by the smallest eigenvalue $\mu_p$: as $\gamma \to 1^-$, $\mu_p/n \to (1-\sqrt\gamma)^2 \to 0$ and the variance term $1/\mu_p$ explodes. Ridge regularization replaces $1/\mu_j$ by $\mu_j/(\mu_j + \lambda)^2$ in each spectral component, which is bounded by $1/(4\lambda)$ regardless of $\mu_j$. The cost is a bias term, but~\eqref{eq:ridge-bias-bound} shows the bias is universally bounded by $\opnorm{\bb}^2$. Choosing $\lambda$ to balance the two leads to a finite-risk estimator even at $\gamma = 1$.

\emph{Optimal $\lambda$ (sketch).} Differentiating $\opnorm{\bb}^2 + \sigma_\varepsilon^2 p/(4\lambda)$ in $\lambda$ does not yield a finite minimizer (the bound is monotone decreasing in $\lambda$ within $\lambda > 0$, but with the actual bias~\eqref{eq:ridge-bias} the trade-off is delicate). Under the isotropic-design assumption $\bX'\bX/n \approx \bI_p$ (so $\mu_j \approx n$ for all $j$), the bias~\eqref{eq:ridge-bias} simplifies to $\lambda^2 \opnorm{\bb}^2/(n+\lambda)^2$ and the variance to $\sigma_\varepsilon^2 p \cdot n/(n+\lambda)^2$; the sum is minimized at $\lambda^* = p\sigma_\varepsilon^2/\opnorm{\bb}^2$, giving total risk $\asymp \opnorm{\bb}^2\sigma_\varepsilon^2 p/(n\opnorm{\bb}^2 + \sigma_\varepsilon^2 p)$. This is bounded uniformly in $\gamma$.

\textbf{Solution to Exercise~\ref{ex:mC3}.} For Rademacher $\varepsilon_i$, $S_n = \sum \varepsilon_i$. The MGF is $\expc[e^{tS_n}] = (\cosh t)^n$. By the Chernoff bound: $\prb(S_n \geqslant s) \leqslant \inf_{t>0}e^{-ts}(\cosh t)^n$. Using $\cosh t \leqslant e^{t^2/2}$ gives $\prb(S_n \geqslant s) \leqslant e^{-s^2/(2n)}$, matching Hoeffding. The bound $\cosh t \leqslant e^{t^2/2}$ is tight at $t = 0$ (both sides equal $1$, derivatives match), so the Hoeffding bound is essentially optimal for Rademacher sums.

\textbf{Solution to Exercise~\ref{ex:mC4}.} (a) $\hat\beta_j = \mathbf{e}_j'\hbb = \mathbf{e}_j'(\bX'\bX)^{-1}\bX'\bY = \beta_j + \mathbf{e}_j'(\bX'\bX)^{-1}\bX'\be = \beta_j + \bc'\be$ where $\bc = \bX(\bX'\bX)^{-1}\mathbf{e}_j$.

(b) Conditionally on $\bX$, $\bc'\be = \sum_i c_i\varepsilon_i$ is a sum of independent sub-Gaussian variables with $\|\bc'\be\|_{\psi_2} \leqslant C\sigma\|\bc\|$. By the sub-Gaussian tail bound:
\[
\prb(|\hat\beta_j - \beta_j| \geqslant t \mid \bX) \leqslant 2\exp\!\left(-\frac{ct^2}{\sigma^2\|\bc\|^2}\right).
\]
Now $\|\bc\|^2 = \bc'\bc = \mathbf{e}_j'(\bX'\bX)^{-1}\bX'\bX(\bX'\bX)^{-1}\mathbf{e}_j = [(\bX'\bX)^{-1}]_{jj}$.

(c) Setting $2\exp(-ct^2/(\sigma^2[(\bX'\bX)^{-1}]_{jj})) = \alpha$ gives $t = \sigma\sqrt{[(\bX'\bX)^{-1}]_{jj}\cdot\frac{1}{c}\ln(2/\alpha)}$. The confidence interval is:
\[
\hat\beta_j \pm \sigma\sqrt{\frac{[(\bX'\bX)^{-1}]_{jj}}{c}\ln\frac{2}{\alpha}}.
\]
\emph{Comparison.} The classical interval has width $\propto t_{\alpha/2,n-p}\sqrt{[(\bX'\bX)^{-1}]_{jj}}$, using $s$ in place of $\sigma$. For large $n-p$, $t_{\alpha/2,n-p} \approx z_{\alpha/2} = \Phi^{-1}(1-\alpha/2) \approx \sqrt{2\ln(2/\alpha)}$ (by the Gaussian tail approximation). So the modern and classical widths agree up to constants. The modern bound has two advantages: (1)~it does not require normality, and (2)~it is a finite-sample statement, not asymptotic. Its disadvantage: it uses the unknown $\sigma$ rather than the data-driven $s$ (though $s$ concentrates around $\sigma$ by Exercise~\ref{ex:mC1}).

\end{document}
